\documentclass{article}
\usepackage{graphicx} % Required for inserting images

\title{Freelancing}
\author{Riccardo Parrino}
\date{May 2025}

\begin{document}

\maketitle

\section{Diritto commerciale}

\subsection{Pillole di diritto commerciale}

L'imprenditore all'art.2082: imprenditore \'e chi esercita progessionalmente un'attitivit\' economica organizzata al fine della prodczione e dello scambio di beni e eservizi

attivit\'a economica: sopportare cost e realizzare ricavi.

Attivit\'a organizzata: presenza di capitale e elavoro.

Attivit\'a professionale: svolta in modo abitale, costante, duraturo, non occasionale anche se stagionale.

Classificazione di imprese:
impresa commerciale: imprenditore non piccolo, non agricola: industriale, trasporto per terra, acqua, mare o aria, attivit\'a bancaria, assicurativa, acquisto e vendita beni.

impresa agricola: \'e imprenditore agricolo chi esercita un'attivit\'a diretta alla coltivazione del fondo.

Per le imprese agricole: no obbligo di tenuta delle scritture contabili, no fallimento, sia piccole che grande imperse.

Piccolo imprenditore: coltivatore diretto del fondo, artifiano, piccoli commercianti. Il piccolo imprenditore non \'e soggetto a pubblicit\'a legale,
non \'e soggetto a fallimento in determinati parametri.

In particoalre, i piccoli imprenditori sono solo gli imprenditori individuali, non i collettivi.

Grande imprenditore: coloro che significativo investimento di capitale, con fallimento ed altre procedure concorsuali.

Sono soggetti all'obbligo delle iscrizione nel registro delle imprese.

Grandi o medie, legalmente non vi \'e differenza.

Classificazione dell'impresa per relazione al soggetto esercente: impresa pubblica e impresa privata.
Le prime sono imprese esercitate fagli enti pubblici, distinguiamo inoltre gli enti pubblici economici dagli enti pubblici non economici.

E le imprese private, cio\'e gestite da privati.


Registro delle imprese:
strumento di pubblicit\'a legale. serve a registrare i fatti e gi atti pi\0u sgnificativi. Si trova presso l'ufficio del registro. Chi h 

Obbligo della registrazione: imprenditori commerciaaali individuali, societ\'a commerciali, mutualistiche, enti pubblici enconomici, societ\' bancarie, assicurative, attivit\'a ausiliarie delle precedente, i consorzi. 

Tenuta informatica delle scrittture contabili: vidimizaione, numerazione e tenuta infoemarica.

Scritture contabili
documenti che contengono la situzione patrimoniale dell'impresa. Rappresentano il modo in cui sta andndo l'azienda. Rendere effeciente l'organizzazione e la gestione dell'imrprea. Tutela gli stessi imprenditori tra creditori e debitori. 
Imprenditore commerciale: ha l'obbligio di tenere il libro giornale, lirbo degli inventari (attivit\'a e passivit\'a dell'imprenditore). 
Altre scritture contabili: libro mastro, libro magazzino, libro cassa, corrsipndenze: lettere, telegrammi, fatture ricevute, fatture spedite, corretta tenuta delle scritture contabili.



Societ\'a in generale
2247 cc: societ\'a contratto tra due o pi\'u persone. Quando cresce una societ\'a  tra le parti.

Contratto consensulae, esercizio in comune di un0attivit\'a economica, conferminto di beni e serviz da parte dei soci. 
Scopo lucrativo: divisione utili tra i soci
Patrimonio sociale: beni conferiti in propriet\'a non appartengono pu\0u ai singoli soci
gesione comune dell'attivit\'a econmica.


Societ\'a di persone: societ\'a semplice, nome collettivo, accomnadita semplice
non sono persone persone giuridiche, scopo di lucro, autonomia patrimonaile imperfetta, prevalenza dell'elemento personale sul patrimonio. 


Societ\'a di capitali:
per azioni, in accomndita per azione
a responsabilit1'a limitata.
sono persone giuridiche, scopo di lucro, autonomia perfetta


Autonomia patrimoniale: 
imperfetta: patrimonio ditta non \'e separato da quello dei soci.
perfetta: patrimonio della socit\'a \'e separato da quello dei soci. 


SAPA e SRL
combinazione tra spa e sas, azionisit accomandatari e accomandanti

SRL
possiede quote di partecipazione e non azioni
capitale minimo di 10.000 euro
trarre mezzi finanziari dalle risorse di un ristretto gruppo di soci. 

srl a socio unico o unipersonale.
iscirzione nel registro delle imprese
autonomia patrimoniale

ditta individuale:
definizione e requisiti di costituzione: iscrizione al registro delle imprese, apertura di una partita IVA, autorizzazione specifiche se necessarie.
caratteristiche principali: responsabilit\'a illimitata, autonomia gestionale, fiscalit\'a semplice in base al regime scelto.
Vantaggi dell'impresa individuale: semplicit\'a burocratica, costi ridotti, flessibilit\'a operativa, possibilit\'a di agevolazioni fiscali, come il regime forfettario.
Svantaggi e rischi: responsabilit\'a patrimoniale illimitata, limitata capacit\'a di acesso al credito, crescita limitata, carico fiscale crescente. 
Se l'azienda individuale cresce notevolmente \'e necessaria ed indicata la trasformazione in un'azienda di capitali, per la separazione del patrimonio personale da quello dell'impresa. Questa operazione richiede procedure formali e approvazione di un notaio.

I diversi modi con la quale \'e possibile avviare un business sono i seguenti, per ordine di complessit\'a:
partita IVA, ditta individuale, srls, srl, Srl (Startup Innovativa), S.P.A.


\subsection{Libero Professionista o Ditta individuale}
Per aprire una ditta individuale \'e necessario completare alcuni pochi passaggi. Per prima cosa scegliere una ragione sociale della ditta,
Il secondo passaggio consiste nel presentare la Comunicazione Unica al Registro delle imprese, presentadosi come imprenditore commerciale o come piccolo imprenditore.
Attraverso questo secondo passaggio, sar\'a possibile adempiere a tutte le formalit\'a necessarie per costituire un'unica impresa: attribuzione Partita Iva, iscrizione al Registro delle Imprese ed eventualmente all'Albo delle Imprese Artigiane,
adempimenti INPS ai fini previdenziali, adempimenti INAIL ai fini assicurativi, e Segnalazione Certficata di Inizio Attivit\'a (SCIA) per lo Sportello Unico delle Attivit\'a Produttive (SUAP).
Sar\'a inoltre necessario avere un acasella di Posta elettronica certificata (PEC) e di un dispositivo per la Firma digitale del titolare d'impresa o del suo delegato.

I costi di apertura di un'azienda individuale sono relativamente bassi e si possono dividere in costi fissi e variabili. 
I costi fissi riguardano l'imposta di bollo, i diritti di segreteria della Camera di Commercio ed il diritto camerale. 
I costi variabili cambiano a seconda dell'attivit\'a esercitata e sono l'Attestazione di inizio attivit\'a e registrazione al SUAP, se l'attivit\'a esercitata \'e soggetta a SCIA.

Per quanto riguarda la tassazione, dipende dal regime fiscale alla quale si decide di aderire, ma supponendo una fase iniziale ed agevolata, \'e consiglabile un regime forfetario.
Secondo questo regime \'e prevista un'imposta sostitutiva al 15\% di tutte le imposte ordinario, applicata alla somma dei ricavi e dei compensi ottenuti dalla ditta all'interno di un anno, moltiplicati per il coefficienti di redditivi\'a, deciso dal codice ATECO della professione. 

A questi costi si aggiungono i contributi INPS, che seguono gli stessi principi e regole che valgono per i liberi professionisti.

Se si rientra nella tipologia di commerciante o artigiano allora si dovr\'a aprire partita iva come ditta individuale.

\subsection{Societ\'a a responsabilit\'a limitata SRL}
Piena personalit\'a  giuridica, i creditori possono rifarsi solo sul patrimonio aziendale. 
Offre flessibilit\'a, scalabilit\'a, 
accesso al credito (banche ed altri istituit di credito concedono pi\'u facilmente prestiti alle srl, considerate un rischio minore), 
successione aziendale (la trasformazione in srl facilita il passaggio della propriet\'a ad altri soci o eredi), tassazione (per pianificazione fiscale).

Nata per facilitare la costituzione di nuove imprese, per piccoli imprenditori, specialmente startup.

I soci di una srls devono essere solo ed esclusivamente persone fisiche. Nella Srl ordinaria i soci possono essere anche eprsono giuridiche.

\subsection{Societ\'a a Responsabilit\'a limitata semplificata SRLs}

\subsection{Societ\'a SRL (Startup)}

\subsection{Societ\'a per azioni S.P.A.}

\section{Diritto tributario}

\section{Ragioneria}

\subsection{Pillole di ragioneria}
Fasi di vita aziendale: organizzazione, gestione e rilevazione

rilevazione: rappresentazione dei dati che emergono dalle op. aziendali
fase tipica e fondamentale dell'azienda
alla base della rilevazione, il sistema informativo aziendale
persone, software e strumenti per la raccolta e l'analisi di queste informazioni

dare informazioni, essere disponibili, per interscambio tra le unit\'a operanti

aggregazione dei dati, quindi informazioni per il suppot delle operazioni e comunicazione con l'esterno

produrre e gestire una notevole quantit\'a di informaizoni, limitare gli errori, smistare le info in tempo reale, raccolta di info con modalit\'a prefissate

si ottimizzano le informationi: sistema dei debiti, il sistema dei crediti, il sistema dei cespiti, il sistema del personale, logistico e del magazzino, finanziario per entrate e euscite, ricerca e sviluppo, marketing ecc

offrire e ricevere informaizoni

progettato per prendere decisioni
reuqisito di tempestitvit\'a e dato da tutte le aree della socità

sottosistema contabile e non contabile
contabile: relative alle op. monetarie
inventari, scritture elementari, contabilit\'a sezionali, contabilit'a generale, contabilit\'a analitico-gestionale, pianificazione e controllo budgetario

sistema non contabile: statistiche, previsioni andamenti futuri, simulazione altrenative direzionali, analisti di mercato o di settore, reporting

rilevazione
rilevazione = scrittura
ha due finalit\'a: consuntive(rappresentare i dati) e di controllo

documenti originari: di prova, di autorizzazione, di controollo
fatrture, contratti, documenti di trasporti, gli assegni, le cambiali

classificazione delle scritture (rilevazioni)
grado di elaborazioni dati: elementari, sezionali, complesse
ordine logico: cronologiche, sistemtiche
obbligatoriet\'a: obbligatorie e facoltative
strumento di rilevazione: contabili, extracontabili

libro cassa e lo scadenziario dei debiti
libro giornale: scrittura complessa, le scritture pi\'u importanti enlla vita della societ\'a
gli stakeholder valutano l'andometno sulla basse delle scritture complesse

scritture cronologiche come il libro magazzino
scritture sistematiche

scritture obbligatorie

scritture contabili: usano ilconto come strumento di scrittura,
extracontabili invece sulla base di grafici e visualizzazioni

conto
definizione: insieme di scritture omogenee realtive ad un determinato oggetto
dare, avere
aprire / accednere un conto
addebitare
accreditare
tenere un conto
chiudere / spegnere un conto

classificazione dei conti:
funzionamento: bilaterali, unilaterali
grado di analisi oggetto: analitici, sintetici (co. ge.)
informaizoni relative all'oggeto sinottici (co. ge.) e descrittivi
forma: sezioni divise/contrapposte (co.ge), sezioni unite/accostate

contabilit\'a sezionali
prima nota / brogliaccio: scrittura elementare riferimento co.ge.(contabilit\'a generale)
mese, giorno, riferimento

contabilit\'a sezionali: settori azienda

contabilit\'a di cassa: entrate/uscite denaro, bolli e assegni, registratori di cassa, reversali/mandati, saldo: sempre positivo
contabilit\'a rapporti con le banche: c/c bancari, consistenza variazioni, estratto conto, operazioni bancarie, saldo: + o -
contabilit\'a fornitori: 
contabilit\'a clienti: documenti vendita e regolazioni
contabilit\'a IVA: 
contabilit\'a magazzino: movimento dei beni
contabilit\'a beni strumentali: registro cespiti ammortizzabili, libro inventari, registro acquisti: storia dei beni
contabilit\'a personale: libro matricoola, libro paga, registro infortuni

Obblighi contabili
adempimneti che le imprese sono tenute ad effettuare per manetenere coerenza tra libro contabile e 
normaitva civilistica
legislazione del lavoro
normativa fiscale
altri obblighi

Normativa civilistica:
art. 2082 c.c. -> imprenditore commerciale
art. 2195 c.c. -> impresa commerciale
art. 2214 c.c. -> imprenditore commerciale deve tenere il libro giornale ed il libro inventari
art. 2220 c.c. -> imp. commerciale deve tenere anche: libri contabili specifici azienda, per 10 anni atti che riguardano l'attivit\'a
art. 2083 c.c. -> piccolo imprenditore: coltivatore diretto, artigiano, lavoratore in proprio

scrittura contabili previsti da normativa a civilistica: 2114 a 2220
libro giornale, libro inventari, scritture contabili richieste dalla natura e dalle dimensioni dell'impresa (es. socit1'a di capitali: libro soci, libro verbali adunanze)

Legislazione del lavoro: libro matricola e libro paga: sostuiti dal libro unico, soggetti a bollatura da parte dell'INAIL
registro degli infortuni

Normativa fiscale: 
DPR 633/72 decreto iva, DPR 600/73 decreto accertamento imposte sui redditi
Decreto IVA: registro acquisti, registro fatture emesse, registro corrispettivi di vendita

regime contabile agevolato per contribuenti minimi

contabilit\'a ordinaria -> societ\'a capitali, individuali, di persone

finalit\'a tenuta e conservzione delle scritture contabili obbligatorie

art. 2219 c.c. contabilit\'a ordinata
art. 2220 c.c. conservazione per 10 anni

contabilit\'a generale
co. ge.
Principale strumento rilevazione: registrare fatti esterni di gestione
metodi di scritture: regole
metodo della partita doppia \'e applicato al sistema del capitale a patrimonio e del risultato economico
(ideato da Amaduzzi)

basato sul dare ed avere
variazioni finanziarie vs variazioni economiche

Partita doppia, conti finanziari ed economici

\section{Contabilit\'a}
\subsection{Pillole di contabilit\'a}
Esistono diverse contabilit\'a per le aziende: Mercantile o commerciale, Industriale o tecnica.

Per mercantile o commerciale, si intende la contabilit\'a che collega l'azienda con il mondo esterno e riguarda l'acquisto di materie prime (uscite) e la vendita di prodotti o servizi (entrate);

Per industriale o tecnica, si intende la contabilit\'a che sta all'interno dell'azienda e riguarda i processi di trasformazione delle materie prime in prodotti finiti.

Da questa distinsione ne segue un'ulteriore, ovvero: la contabilit\'a generale (Co.Ge) e la contabilit\'a industriale (anche detta contabilit\'a analitica, Co.An).
La prima registra tutti i fatti amministrativi intercorsi tra l'azienda e l'ammbiente esterno. La seconda registra solo i fatti di gestione interna.

Nella Pubblica Amministrazione la metodologia viene detta della Contabilit\'a Finanziaria, e fino a pochi anni fa era della partita semplice. Da pochi anni \'e stato introdotto il metodo della partita doppia, con la denominazione contabilit\'a economica.

Ragioneria analizzare la contabilit\'a 
cnotabilit\'a: tutte le operazioni che un 'azienda svolge nella sua vita
metodo della partita doppia per realizzare ci\'ob
conti raccolti nel libro mastro. In due sezione, dare ed avere
ogni fatto di gestione va osservato sconod due aspetti, aspetto originario ed aspetto derivato
il tatolae degli importi "dare" deve essere uguale al totale importi "avere"

esistono diversi libri, libro giornale, libro mastro e libro inventari.
libro giornale: le operazioni fatte ogni giorno

libro inventari: obbli. all'inizio dell'impresa, insieme di tutte le passivit\'a dell'azienda

classificazioni di operazioni:
variazioni finanziarie e variazioni ECOnomiche
finanziarie: positive(aumento di denaro) e negative(diminuzione di denaro)

var. economiche: positive(ricave, aumento di capitale), negative(costi, diminuzione di capitale)

conto: + (dare, vfp, ven), - (avere, vfn, vep)

iva

Operazioni di acquisto
acquisto di beni
acquisto di servizi
acquisto d immobilizzazioni
acquisto sui mercati esteri

ogni operazoini va rilevata
correttta valorizzazione dell'immobile
momento di rilevazione
valori originati: costo del bene, debito di regolamento, iva a credito
fasi rilevazione: liquidazione, pagamento

Esempio di acquisto di beni
acquisto di beni
mat.prime c/acquisti +  iva a credito: dare 10.000 + 2.200
avere 12.200

\section{Avvio attivit\'a}
\subsection{Partita IVA}
\subsection{Apertura partita iva}
Per aprire partita iva \'e necessario inviare il modello AA9/12 o telematicamente o fisicamente all' Agenzia delle entrate. 

All'interno di questo modello \'e necessario inserire il proprio nome e cognome se si tratta di lavoratore autonomo, della denominazione dell'azienda se si tratta di un'impresa. Successivamente \'e necessario inserire il codice ATECO della propria attivit\'a, il tipo di regime fiscale (ordinario, ordinario semplificato o forfettario, supponendo che si soddisfino tutti i requisiti), il luogo in cui si svolge l'attivit\'a ed eventuali collaboratori come anche il prorpio commercialista.

Per quanto riguarda il luogo di svolgimento della propria attivit\'a \'e necessario indicare un luogo in cui si \'e facilmente reperibili ed \'e apposta eventuale targhetta di riconoscimento. Se si \'e in affitto, \'e necessario informare il proprietario di casa che si intende usare l'abitazione come ufficio.

Oltre a questi adempimenti in talune situazioni \'e necessario iscriversi all'INPS per gestione separata o artigiani/commercianti, all'INAIL in caso di rischio per la salute, alla Camera di Commercio (per artigiani e commercianti) e SCIA al Comune in caso di  attivit\'a soggetta a comunicazione di inizio.

Successivamente a questi passaggi \'e necessario avere una propria PEC (Posta Elettronica Certificata) e se non si rientra nel regime forfettario \'e obbligatoria la fatturazione elettronica tramite SDI. 

In genere i tempi di apertura sono immediati ed i costi sono ridotti ai soli servizi professionali (come il commercialista), i contributi INPS ed eventuali diritti camerali.

\subsection{Chiusura partita iva}
Inviare il modello AA9/12 di chiusura, in maniera analoga a quanto fatto per l'apertura. Non ci devono pi\'u essere fatture da emettere. Per almeno un anno, bisogno mantenere i rapporti con il proprio consulente, per la dichiarazione dei redditi e eventuali F24 da pagare per l'anno precedente.

\subsection{Codice ateco}
Il codice ATECO (acronimo di ATtivit\'a ECOnomiche) \'e una classificazione alfanumerica usata in Italia per identificare il tipo di attivit\'a economica svolta da un'impresa o da un lavoratore autonomo. Il codice ATECO \'e fondamentale per l'apertura della partita IVA, la determinazione del regime fiscale, contributivo e per fini statistici. 

Il codice ATECO ha una struttura gerarchica XX.YY.ZZ.W, in cui XX individua la sezione, YY individua la divisione, ZZ il gruppo e W la Classe/Voce, anche se quest'ultimo non \'e sempre presente e serve a dare un ulteriore livello di dettaglio. Un esempio pratico: Codice ATECO: 74.10.21 significa 74 Altre attivit\'a professionali, scientifiche e tecniche, 10 si intende Attivit\'a di design specializzate e 21 Design grafico, in particolare per il web. Il codice ATECO si usa per l'apertura della partita IVA, per il calcolo dei contributi INPS, per la scelta del regime forfettario,, per statistiche ISTAT ed in alcuni casi per autorizzazione e licenze.

Per attivit\'a legate allo sviluppo software o legate all'informatica, la divisione di riferimento sono la 62 e la 63. 
Per una spiegazione dettagliata e precisa, fare ricerche al seguente link: https://www.istat.it/classificazione/classificazione-delle-attivita-economiche-ateco


\section{Regime fiscale}

\subsection{Regime fiscale ordinario}
\subsubsection{Requisiti}
Il regime ordinario \'e obbligatorio per le societ\'a di capitali e 
facoltativo per le societ\'a di persone e ditte individuali che nell'anno precedente non abbiano conseguito ricavi superiori a: 
500.000 euro nel caso di prestazione di servizi e 800.000 euro negli altri casi. Superati qeusti limiti il regime ordinario diventaobbligatorio anche per le societ\'a di persone.

\subsubsection{Fiscalit\'a}
Le imposte previste nel regime ordinario sono: IRPEF (per le persone fisiche e le ditte individuali), 
IRES (per le societ\'a), IRAP (per le societ\'a), IVA (su tutte le fatture attive e passive).

Per il calcolo delle tasse nel regime ordinario vale il principio di competenza, il sistema di proporzionalit\'a ed i costi sono deducibili.

Per quest'ultima voce si intende il fatto che nel regime ordinario \'e possibile dedurre diversi costi per la gestione dell'attivit\'a, concorrendo all'abbassamento del reddito imponibile.

Per persone fisiche \'e previsto il pagamento di:
\begin{itemize}
    \item IRPEF (su reddito imponibile)
    \item IVA
    \item addizionali comunali
    \item ritenute d'acconto
    \item Contributi previdenziali presso la propria cassa
    \item se si hanno dei dipendenti, bisogna adempiere ai vari obblighi verso INAIL ed INPS dipendenti.
\end{itemize}

Per societ\'a (sia di persone che di capitali) sono previsti:
\begin{itemize}
    \item IRES (su reddito imponibile)
    \item IRAP (su reddito imponibile)
    \item IVA
    \item addizionali comunali
    \item ritenute d'acconto
\end{itemize}

\subsubsection{Fiscalit\'a: Quali costi sono deducibili?}
Si distingue in spese detraibili e spese deducibili: 
le spese deducibili riducono il reddito imponibile (ovvero la base su cui viene calcolata l'imposta),
mentre le spese detraibili riducono direttamente l'imposta dovuta, creando uno sconto sull'importo finale che si deve versare.

\subsubsection{Costi deducibili e costi detraibili}
Le spese deducibili variabo in base  alla forma giuridica utilizzata ed al regime fiscale adottato  (libero professionista, ditta individuale o societ\'a, semplificato o ordinario). 

Tuttavia, la regola generale stabilisce che un costo \'e deducibile se rispetta il principio di inerenza.

Un caso \'e inerente all'attivit\'a d'impresa se rispetta i criteri di effettivit\'a, inerenza e cnogruit\'a. 
Effettivit\'a: il costo \'e effettivamente stato sostenuto; inerenza: il costo \'e strettamente connesso all'attivit\'a professionale o imprenditoriale; congruit\'a: i costi sono proporzionati all'attivit\'a svolta.

Non tutti i costi sono deducibili all 100\%, alcuni di questi lo sono solo parzialmente.

I seguenti costi sono deducibili al 100\%: materie prime e merci, materiale di concelleria, beni strumentali fino a 516,36 euro, investimenti in pubblicit\'a e marketing, consulenze professionali per supporto legale, contabile o tecnico, contributi previdenziali (verso INPS o altri enti previdenziali), utenze aziendali come acqua, luce, gas esclusivi per l'attivit\'a.

Tuttavia, la deducibilit\'a pu\'o variare in base al regime fiacale adottato.

Per costi parzialmente deducibili si intendo invece i csoti che non sono interamente deducibili dal reddito imponibile. Alcuni esempi di questi sono le spese di rappresentanza, alcuni benefit ai dipendenti, le autovetture e i veicoli, le spese telefoniche.

Ricapitolando, vi \'e una sostanziale differenza tra costo deducibile e costo detraibile: i costi deducibili sono i costi che riducono il reddito imponibile ed i costi detraibili sono i costi che riducono l'imposta dovuta.

\subsubsection{Contabilit\'a}
Per quanto riguarda la contabilit\'a esistono diversi adempimenti obbligatori, ovvero: la fatturazione elettronica, l'applicazione dell'IVA in fattura, la dichiarazione IVA (presso l'Agenzia delle Entrate), il versamento dell'IVA (a cadenza mensile o trimestrale)
e la compilazione di libri e registri contabili.

\subsubsection{Contabilit\'a: Quali libri compilare}
Ai fini fiscali \'e prevista la compilazione dei libri: libro giornale, 
libro degli inventari, 
scritture ausiliarie (libro mastro e scritture di magazzino), 
registro dei beni ammortizzabili, 
registri previsti dalla normativa Iva (registro delle fatture emesse, registro dei corrispettivi, registro degli acquisti), 
libri sociali obbligatori previsti ai fini civilistici.
Per ulteriori precisazioni, si rimanda al commercialista.

Quindi ricapitolando, in regime ordinario \'e necessario seguire gli adempimenti:
\begin{itemize}
    \item registrazione fatture attive e passive
    \item fatturazione elettronica
    \item applicazione dell'IVA in fattura
    \item dichiarazione IVA annuale
    \item versamento IVA
    \item compilazione libri e registri contabili
    \item redazione di bilancio d'esercizio (per le societ\'a)
    \item dichiarazione dei redditi
    \item dichiarazione IRAP
    \item spesometro / esterometro (per operazioni estere)
    \item comunicazioni INTRASTAT
\end{itemize}

Per quanto riguarda gli adempimenti ai fini contabili, \'e necessario tenere:
\begin{itemize}
    \item libro giornale
    \item libro degli inventari
    \item scritture ausiliarie (conti mastro, scritture di magazzino)
    \item registro dei beni ammortizzabili
    \item registri previsti dalla normativa IVA
    \item libri sociali obbligatori previsiti ai fini civilistici
\end{itemize}

\subsection{Regime fiscale ordinario semplificato}
\subsubsection{Requisiti}
Vi possono irentrare tutte le imprese individuali e le societ\'a di persone se i loro ricavi nell'arco di un anno solare 
non superano i seguenti limiti: 500.000 euro per le prestazioni di servizi e 800.000 euro per tutte le altre attivit\'a.
Per il professionisti non vige alcun limite.

Si accede al regime semplificato indicando un volume d'affari che non supera le soglie di ricavi inferiore a quelle indicate al momento dell'apertura della partita IVA.

\subsubsection{Fiscalit\'a}
La tassazione segue le stesse regole previste il regime ordinario. Per le persone fisiche, la tassazione viene determinata in modo progressivo, basandosi sulle aliquote IRPEF.

Ricordando, per persone fisiche, le imposte sono:
\begin{itemize}
    \item IRPEF (su reddito imponibile)
    \item Contributi Previdenziali
    \item IVA
\end{itemize}

Per societ\'a (di capitali o persone), si devono pagare:
\begin{itemize}
    \item IRES
    \item IRAP
    \item IVA
\end{itemize}

\subsubsection{Contabilit\'a}
A differenza del regime ordinario, nel semplifcato non si ha l'obbligo di alcune scritture contabili e non si ha l'obbligo di redigere il bilancio a fine anno.

Gli obblighi di contabilit\'a in questo regime prevedono: la tenuta del registro dei beni ammortizzabili, la tenute del registro IVA, la tenuta del registro degli incassi dei pagamenti e la tenuta del libro unico del lavoro. 

Anche nel regime semplificato (come in tutti gli altri regimi) \'e obbligatorio la fattura elettronica.

In conclusione, si devono tenere:
\begin{itemize}
    \item il registro dei beni ammortizzabili
    \item registro IVA
    \item registro degli incassi dei pagamenti
    \item libro unico del lavoro
    \item fatturazione elettronica
\end{itemize}

\subsection{Regime fiscale forfetario}
Come riportato sul sito dell'agenzia delle entrate, "il regime forfetario prevede rilevanti semplificazioni ai fini IVA e ai fini contabili, e consente, altresì, la determinazione forfetaria del reddito da assoggettare a un’unica imposta in sostituzione di quelle ordinariamente previste, nonché di accedere ad un regime contributivo opzionale per le imprese."
Inoltre, \'e possibile accedere al regime forfetario senza vincoli di tempo ne di et\'a, ma soltanto soddisfando i requisiti previsti.
Il regime forfetario \'e stato introdotto dalla legge di stabilit\'a 2015 e rappresenta il regime naturale delle persone fisiche che esercitano un'attivit\'a di impresa, arte o professione in forma individuale.

Per quanto riguarda reddito e tassazione, per chi aderisce al regime forfetario \'e prevista un'unica imposta sostitutiva delle imposte sui redditi, delle addizionali regionali e comunali e dell'IRAP, nella misura del 15\%. 
L'aliquota viene applicata sull'imponibile determinato dall'ammontare dei ricavi o dei compensi percepiti moltiplicato per il coefficiente di redditivit\'a, diversificato a seconda del codice ATECO che contraddistingue l'attivit\'a esercitata.

Per i codice ATECO 62-63 \'e previsto un coefficiente di redditivit\'a del 67\%.

Per chi aderisce al regime forfetario, per i primi cinque anni di attivit\'a si potrebbe avvalere anche di un ulteriore detassazione, un'imposta sostitutiva del 5\%, visti alcuni requisiti da soddisfare.

Il regime forfetario \'e di pi\'u semplice gestione rispetto al regime ordinario od ordinario semplificato, sia per fini IVA, sia per fini sulle imposte dirette. 

Per quanto riguarda i fini Iva, non deve essere addebitata l'Iva in fattura e non deve essere detratta l'iva sugli acquisti. Non si liquida l'imposta, non si versa e non si \'e obbligati a presentare la dichiarazione e la comunicazione annuale Iva. 
Non \'e necessario comunicare all'Ageniza delle entrate le operazioni rilevanti ai fini Iva (ovvero lo spesometro) né quelle effettuate nei confronti di operatori economici aventi sede, residenza o domicilio in Paesi cosiddetti black list. Chi applica il regime forfetario, inoltre, non ha l'obbligo di registrare i corrispettivi, le fatture emesse e le ricevute.

Le semplificazioni ai fini delle imposte sul reddito riguardano l'esonero dagli obblighi di registrazione e tenuta delle scritture contabili, ma mantenere e cnoservare i registri previsti da disposizioni diverse da quelle tributarie.
Non operano le ritenute alla fonte, non si subiscono le ritenute, non si applicano le ritenute alla fonte. 

Ai fini Iva, si ha l'obbligo di numerare e conservare le fatture di acquisto e le bollette doganali, certificare i corrispettivi, 
integrare le fatture per le operazioni di cui risultatno debitori di imposta con l'indicazinoe dell'aliquota e della relativa imposta, da versare entro il giorno 16 del mese successivo a quello di effettuazione delle operazioni, senza diritto alla detrazione dell'imposta relativa.

La legge di stabilit\'a, insieme alla legge di bilancio, costituisce la manovra di finanza pubblica per il triennio di riferimento
e rappresenta lo strumento principale di attuazione degli obiettivi programmatici definiti con la Decisione di finanza pubblica. (fonte: "https://www.rgs.mef.gov.it/VERSIONE-I/attivita_istituzionali/formazione_e_gestione_del_bilancio/bilancio_di_previsione/bilancio_finanziario/BF-2016_2018/LdS2016/index.html#:~:text=La\%20legge\%20di\%20stabilit\%C3\%A0\%2C\%20insieme,la\%20Decisione\%20di\%20finanza\%20pubblica.")

\subsubsection{Adempimenti regime fiscale forfetario}
Per quanto riguarda gli adempimenti generali del regime fiscale forfetario sono previsti i seguenti: apertura della Partita IVA, fatturazione, registrazione delle operazioni, adempimenti fiscali, calcolo del reddito imponibile e i contributi INPS.

Per l'apertuea della partita iva \'e necessaria la scelta del codice ateco e la comunicazione dell'inizio attivit\'a presso l'Agenzia delle Entrate.

Per la fatturazione \'e obbligatoria l'emissione di fatture elettroniche, in cui vanno inseriti tutt gli elementi della fattura ed in cui bisogno applicare la marca da bollo da 2 euro per fatture superiori a 77,47 euro.

Per la Registrazione delle Operazioni, non \'e obbligatorio tenere registri IVA o libri contabili, ma \'e utile tenere traccia degli incassi e dei pagamenti per una gestione accurata e per future verifiche fiscali. 
In ogni caso \'e necessario conservare tutte le fatture emesse e ricevute.

Per gli adempimenti fiscali, i forfettari no applicano l'IVA sulle fatture e non la detraggono sugli acquisti. E' prevista l'imposta sostitutiva per il pagamento delle imposte e dei contributi previdenziali presso le proprie casse di appartenenza.

Per le dichiarazioni e le comunicazioni, sono previste le dichiarazione dei redditi e la certificazione unica se si hanno dipendenti o collaboratori.

\subsubsection{Scadenze regime fiscale forfetario}

\begin{comment}
    tema imposte
    tema contributi
\end{comment}

I principali adempimenti di chi aderisce al regime forfetario sono il versamento delle imposte, il versamento dei contributi previdenziali, il versamento del diritto comerale annuale, 
invio delle fatture elettroniche e la dichiarazione dei redditi.

Per quanto riguarda il versamento delle imposte, ci sono due principali scadenze: il 30 giugno ed il 30 novembre.
Il 30 giugno si versa il saldo dell'imposta sostitutiva riferita all'anno precedente e il 50\% dell'acconto dell'imposta sostitutiva dell'anno in corso. I versamenti del 30 giugno possono essere rateizzati fino a 6 rati, con cadenza mensile.

L'altra scadenza \'e invece quella del 30 novembre, procedendo al versamento del restante 50\% dell'acconto per le imposte dell'anno in corso. Quest'ultimo versamento non pu\'o essere rateizzato.

Per chi apre partita IVA, durante il primo anno non dovr\'a pagare alcuna imposta sostitutiva durante il primo anno, 
ma durante il secondo anno verser\'a l'ammontare del primo anno e l'acconto del secondo anno.

I contribuenti forfetari pagano un'imposta sostitutiva di IRPEF e addizionali regionali/comunali pari al 15\%, su un imponibile determinato sulla base del totale dei ricavi e dei compensi incassati (vige il principio di cassa) applicando la percentuale di redditivit\'a,
determinata sulla base del codice ateco, e deducendone i contributi previdenziali versati nel periodo.

Per principio di cassa si intende il sistema di registrazione dei costi e dei ricavi all'interno del bilancio secondo la quale vengono registrati soltanto i costi ed i ricavi effettivamente pagati o incassati, non soltanto emessi per fattura come avviene per il principio di competenza.
Ripetendo, con il principio di cassa vengono imputate all'interno del bilancio le sole registrazioni che hanno avuto un movimentato di denaro. Esemplificando, una fattura emessa il 31 dicembre, pagata nell'anno successivo, non sar\'a registrata nel bilancio annuale.

Questo sistema ha il vantaggio di dover pagare le imposte solo dopo aver incassato i pagamenti relativi alle fatture emesse.

Oltre al versamento dell'imposta sostitutiva \'e prevista anche l'imposta di bollo sulle fatture elettroniche, ovvero:
per ogni fattura emessa con importo superiore a 77,47 euro, per la cifra di 2 euro per ciascuna fattura. Le scadenze previste sono le seguenti:
per le fatture emesse nel primo, secondo e terzo trimestre, entro il 30 novembre, mentre per le fatture trasmesse nel quarto trimestre (quindi, ottobre-dicembre) entro il 28 febbraio. 

Oltre alle imposte, il soggetto a partita IVA forfetaria \'e tenuto al versamento dei contributi previdenziali.
L'aspetto contributivo dipende dall'attivit\'a svolta dal contribuente.
L'ammontare dovuto ed i termini in cui pagarlo \'e determinato dall'attivit\'a del contribuente.

Innanzi tutto bisogna distinguere in contribuenti che hanno una propria cassa previdenziale ed in coloro che non ce l'hanno e devono quindi affidarsi alla Gestione Separata INPS.
Ed inoltre, tra coloro che non hanno una propria cassa previdenziali, gli artigiani ed i commercianti dai restanti contribuenti. 

Per chi ha una propria cassa previdenziale, come gli avvocati, gli ingegneri e gli architetti ecc., le modalit\'a e le scadenze sono definite dalla cassa previdenziale stessa. 

Per chi invece non ha una propria cassa previdenziale deve iscriversi alla Gestione Separa INPS, pagando ogni anno un versamento in misura percentuale sul reddito forfetizzato di ogni anno.
Le aliquote variano tra il 26,07\% per chi non ha alcuna altra copertura previdenziale ed il 24,00\% per chi \'e pensionato o \'e gi\'a iscritto ad un'altra forma di previdenza obbligatoria.

Le scadenze previste per questi versamenti sono le stesse per quelle previste dall'imposta sostitutiva, 30 giugno e 30 novembre.

Come si diceva, vi \'e un'ulteriore distinsione per Artigiani e Commercianti, secondo la quale \'e prevista un'imposta di reddito detta "minimale" di 18.415 euro, da versare entro l'anno, indipendentemente dal reddito prodotto.
In altre parole, anche se in un determinato anno non si produce alcune reddito, si \'e comunque tenuti a pagare alla Gestione separata INPS, un contributo di 18.415 euro.

Se il proprio reddito, invece, supera l'importo minimale, allora andranno versati ulteriori contributi pari al 24\% (da verificare ogni anno) sul reddito eccedente.

Le scadenze sono articolate in quattro date, ovvero: 16 maggio, 20 agosto, 16 novembre, 16 febbraio (dell'anno successivo).

Per le attivit\'a che devono essere registrate alla camera di commercio (come gli artigiani ed i commercianti) \'e previsto anche il versamento del diritto camerale annuale alla Camera di Commercio. L'importo oscilla tra 44 e 53 euro. L'import deve essere versato entro il 30 giugno di ogni anno.

Dal primo gennaio 2024 anche i contribuenti in regime forfetario devono emettere fatture elettroniche, per beni (al momento della consegna o spedizione, incasso del corrispettivo se precedente la consegna o la spedizione)
e servizi (al momento dell'incasso del corrispettivo). In caso di fatturazione differita di beni o servizi, la fattura deve essere emessa e trasmessa entro il giorno 15 del mese successivo.

Infine, per chi applica il regime forfettario \'e prevista la repsentazione della dischiarazione dei redditi annuale, prevista entro il 30 settembre di ogni anno. 

\subsubsection{Come si pagano le imposte ed i contributi}
Sia le imposte che i contributi possono essere pagate tramite modello F24. 
Il modello F24 pu\'o essere pagato online tramite home banking, tramite commercialista o tramite banca o posta.

In alternativa al modello F24 esistono comunque diversi servizi telematici come FiscoOnline ed Entratel. Spesso comunque, affidandosi ad un commercialista, questo si occupa della trasmissione e del pagamento di contributi e imposte.

\subsection{Riassunto finale}

In generale, si possono riassumere i requisiti, gli adempimenti fiscali e contabili, e le normative, in base al regime fiscale ed alla forma societaria. Di seguito le possibili alternative:

\subsubsection{Libera professione o Ditta Individuale in regime ordinario}
Adempimenti fiscali:
\begin{itemize}
    \item IRPEF (su reddito imponibile)
    \item IVA
    \item addizionali comunali
    \item ritenute d'acconto
    \item Contributi previdenziali presso la propria cassa
    \item se si hanno dei dipendenti, bisogna adempiere ai vari obblighi verso INAIL ed INPS dipendenti.
\end{itemize}

Adempimenti contabili:
\begin{itemize}
    \item registrazione fatture attive e passive
    \item fatturazione elettronica
    \item applicazione dell'IVA in fattura
    \item dichiarazione IVA annuale
    \item versamento IVA
    \item dichiarazione dei redditi
    \item dichiarazione IRAP
    \item spesometro / esterometro (per operazioni estere)
    \item comunicazioni INTRASTAT
    \item libro giornale
    \item libro degli inventari
    \item scritture ausiliarie (conti mastro, scritture di magazzino)
    \item registro dei beni ammortizzabili
    \item registri previsti dalla normativa IVA
    \item libri sociali obbligatori previsiti ai fini civilistici
\end{itemize}


\subsubsection{Libera professione o Ditta Individuale in regime ordinario semplificato}
Adempimenti fiscali:
\begin{itemize}
    \item IRPEF (su reddito imponibile)
    \item IVA
    \item addizionali comunali
    \item ritenute d'acconto
    \item Contributi previdenziali presso la propria cassa
    \item se si hanno dei dipendenti, bisogna adempiere ai vari obblighi verso INAIL ed INPS dipendenti.
\end{itemize}

Adempimenti contabili:
\begin{itemize}
    \item il registro dei beni ammortizzabili
    \item registro IVA
    \item registro degli incassi dei pagamenti
    \item libro unico del lavoro
    \item fatturazione elettronica
\end{itemize}


\subsubsection{Libera Professione o Ditta Individuale in regime forfetario}
Adempimenti fiscali:
\begin{itemize}
    \item Imposta sostitutiva su reddito imponibile
    \item Contributi INPS
\end{itemize}

Adempimenti contabili:
\begin{itemize}
    \item fatturazione elettronica
    \item dichiarazione dei redditi
    \item certificazione unica se si hanno dipendenti
    \item registro degli incassi e pagamenti (consigliato, ma non obbligatorio)
\end{itemize}



\subsubsection{SRL in regime ordinario}
Adempimenti fiscali:
\begin{itemize}
    \item IRES (su reddito imponibile)
    \item IRAP (su reddito imponibile)
    \item IVA
    \item addizionali comunali
    \item ritenute d'acconto
    \item registrazione fatture attive e passive
    \item fatturazione elettronica
    \item applicazione dell'IVA in fattura
    \item dichiarazione IVA annuale
    \item versamento IVA
\end{itemize}
Adempimenti contabili:
\begin{itemize}
    \item libro giornale
    \item libro degli inventari
    \item scritture ausiliarie (conti mastro, scritture di magazzino)
    \item registro dei beni ammortizzabili
    \item registri previsti dalla normativa IVA
    \item libri sociali obbligatori previsiti ai fini civilistici
    \item redazione di bilancio d'esercizio
    \item dichiarazione dei redditi
    \item dichiarazione IRAP
    \item spesometro / esterometro (per operazioni estere)
    \item comunicazioni INTRASTAT
\end{itemize}



\subsubsection{SRL in regime ordinario semplificato}
Adempimenti fiscali:
\begin{itemize}
    \item IRES (su reddito imponibile)
    \item IRAP (su reddito imponibile)
    \item IVA
    \item addizionali comunali
    \item ritenute d'acconto
\end{itemize}

Adempimenti contabili:
\begin{itemize}
    \item il registro dei beni ammortizzabili
    \item registro IVA
    \item registro degli incassi dei pagamenti
    \item libro unico del lavoro
    \item fatturazione elettronica
\end{itemize}



\subsubsection{SRLs in regime ordinario}
Adempimenti fiscali:
\begin{itemize}
    \item IRES (su reddito imponibile)
    \item IRAP (su reddito imponibile)
    \item IVA
    \item addizionali comunali
    \item ritenute d'acconto
\end{itemize}



\subsubsection{SRLs in regime ordinario semplificato}
Adempimenti fiscali:
\begin{itemize}
    \item IRES (su reddito imponibile)
    \item IRAP (su reddito imponibile)
    \item IVA
    \item addizionali comunali
    \item ritenute d'acconto
\end{itemize}



\section{Simulazione}

\subsection{Libero professionista a Regime fiscale forfetario}

\subsection{Ditta individuale a Regime fiscale forfetario}

\subsubsection{Reddito pari a 10.000 euro}
\begin{itemize}
    \item Reddito imponibile: 10.000 * 67\% = 6.700 euro
    \item Imposta sostitutiva (primi 5 anni): 6.700 * 5\% = 335 euro
    \item Contributi INPS (Gestione Artigiani e Commercianti): 2.320,32 euro (con riduzione al 50\%), 
    \item Costo commercialista: 60 euro * 12 mesi = 720 euro
    \item totale utile: (10.000 - 335 - 2320,32 - 720) euro = 6624,68
    \item percentuale costi in tasse e spese: 33,75\%
\end{itemize}

\subsubsection{Reddito pari a 20.000 euro}
\begin{itemize}
    \item Reddito imponibile: 20.000 * 0,67 = 13.400 euro,
    \item Imposta sostitutiva (primi 5 anni): 13.400 * 5\% = 670 euro,
    \item Contributi INPS (Gestione Artigiani e Commercianti): 2.514,32 euro = 194,00 (variabili con riduzione al 50\%) + 2.320,32 euro (con riduzione al 50\%),
    \item Costo commercialista: 60 euro * 12 mesi = 720 euro
    \item totale utile: ( 20.000 - 670 - 2.514,32 - 720 ) euro = 16.095,68 euro
    \item percentuale costi in tasse e spese: 19,52\%
\end{itemize}

\subsubsection{Reddito pari a 30.000 euro}
\begin{itemize}
    \item Reddito imponibile: 30.000 * 0,67 = 20.100 euro
    \item Imposta sostitutiva (primi 5 anni): 20.100 * 5\% = 1.005 euro
    \item Contributi INPS (Gestione Artigiani e Commercianti): 2.514,32 euro = 194,00 (variabili con riduzione al 50\%) + 2.320,32 euro (con riduzione al 50\%),
    \item Costo commercialista: 60 euro * 12 mesi = 720 euro
    \item totale utile: ( 30.000 - 1.005 - 670 - 2.514,32 - 720 ) euro = 25.090,68 euro
    \item percentuale costi in tasse e spese: 16,36\%
\end{itemize}

\subsubsection{Reddito pari a 40.000 euro}
\begin{itemize}
    \item Reddito imponibile: 40.000 * 0,67 = 26.800 euro
    \item Imposta sostitutiva (primi 5 anni): 26.800 * 5\% = 1.340 euro
    \item Contributi INPS (Gestione Artigiani e Commercianti): 3.368,45 euro = 1048,13 (variabili con riduzione al 50\%) + 2.320,32 euro (con riduzione al 50\%),
    \item Costo commercialista: 60 euro * 12 mesi = 720 euro
    \item totale utile: ( 40.000 - 1340 - 670 - 3368,45 - 720 ) euro = 35.241,55 euro
    \item percentuale costi in tasse e spese: 15,24\%
\end{itemize}

\subsubsection{Reddito pari a 50.000 euro}
\begin{itemize}
    \item Reddito imponibile: 50.000 * 0,67 = 33.500 euro
    \item Imposta sostitutiva (primi 5 anni): 33.500 * 5\% = 1.675 euro
    \item Contributi INPS (Gestione Artigiani e Commercianti): 4149,59 euro = 1829,26 (variabili con riduzione al 50\%) + 2.320,32 euro (con riduzione al 50\%),
    \item Costo commercialista: 60 euro * 12 mesi = 720 euro
    \item totale utile: ( 50.000 - 1.675 - 670 - 3368,45 - 720 ) euro = 43.566,55 euro
    \item percentuale costi in tasse e spese: 12,86\%
\end{itemize}

\subsubsection{Reddito pari a 60.000 euro}
\begin{itemize}
    \item Reddito imponibile: 60.000 * 0,67 = 40.200 euro
    \item Imposta sostitutiva (primi 5 anni): 40.200 * 5\% = 2.010 euro
    \item Contributi INPS (Gestione Artigiani e Commercianti): 4969,67 euro = 2649,35 euro (variabili con riduzione al 50\%) + 2.320,32 euro (con riduzione al 50\%),
    \item Costo commercialista: 60 euro * 12 mesi = 720 euro
    \item totale utile: ( 60.000 - 2.010 - 670 - 4969,67 - 720 ) euro = 51.630,33 euro
    \item percentuale costi in tasse e spese: 13,94\%
\end{itemize}

\subsubsection{Reddito pari a 70.000 euro}
\begin{itemize}
    \item Reddito imponibile: 70.000 * 0,67 = 46.900 euro
    \item Imposta sostitutiva (primi 5 anni): 40.200 * 5\% = 2.345 euro
    \item Contributi INPS (Gestione Artigiani e Commercianti): 5789,75 euro = 3469,42 euro (variabili con riduzione al 50\%) + 2.320,32 euro (con riduzione al 50\%),
    \item Costo commercialista: 60 euro * 12 mesi = 720 euro
    \item totale utile: ( 70.000 - 2345 - 670 - 5.789,75 - 720 ) euro = 60.475,25 euro
    \item percentuale costi in tasse e spese: 13,60\%
\end{itemize}

\subsubsection{Reddito pari a 80.000 euro}
\begin{itemize}
    \item Reddito imponibile: 80.000 * 0,67 = 53.600 euro
    \item Imposta sostitutiva (primi 5 anni): 53.600 * 5\% = 2.680 euro
    \item Contributi INPS (Gestione Artigiani e Commercianti): 6.609,82 euro = 4289,50 euro (variabili con riduzione al 50\%) + 2.320,32 euro (con riduzione al 50\%),
    \item Costo commercialista: 60 euro * 12 mesi = 720 euro
    \item totale utile: ( 80.000 - 2.680 - 670 - 6.609,82 - 720 ) euro = 69.320,18 euro
    \item percentuale costi in tasse e spese: 13,34\%
\end{itemize}

\subsection{Ditta individuale a Regime fiscale ordinario semplificato}

\subsection{Ditta individuale a Regime fiscale ordinario}
\subsubsection{Reddito pari a 10.000 euro}
\begin{itemize}
    \item IRPEF 
\end{itemize}

\subsection{SRL a regime fiscale ordinario semplificato}

\subsection{SRL a regime fiscale ordinario}

\section{Fatturazione}

\subsection{Regole della fatturazione elettronica}
A partire dal 2024, tutti i contribuenti in regime forfettario sono obbligati ad emettere fattura elettronica rispettando gli element richiesti dall'art. 21, D.P.R. n. 633/1972, con le seguenti particolarit\'a:
\begin{itemize}
    \item la fattura emessa non deve recare l'addebito dell'imposta;
    \item deve essere assolta l'imposta di bollo di 2 euro, da apporre sull'originale della fattura, qualora la fattua sia di import superiore a 77,47 euro, il cui riaddebito al cliente \'e una facolt\'a
    \item il reverse charge non trova appliczione da parte dei soggetti che operano nel regime forfetario.
    \item i ricavi e i compensi dei contribuenti in regime forfetario non sono assoggettati a ritenuta d'acconto da parte del sostituto d'imposta. Quindi costoro, in fattura, possono indicare che le somme non sono soggette a ritenuta d'acconto (es.: agenti di commercio e professionisti)
\end{itemize}

Si ricorda che l'obbligo generalizzato di emissione della fattura elettronica \'e stato istituito dall'art. 18, D.L. 30 aprile 2022, n. 36, disponendo l'obbligo di emettere fattura elettronica nel rispetto delle seguenti tempisitiche:
\begin{itemize}
    \item a partire dal primo luglio 2022, in capo ai forfettari che, nell'anno precedente, avevano percepito ricavi o compensi superiori a 25.000 euro, ragguagliati ad anno;
    \item a decorrere dal primo gennaio 20204, per tutti i soggetti che adottano il forfetario.
\end{itemize}

\subsection{Elementi da indicare in fattura}
Elementi da indicare in fattura: in base ai chiarimenti forniti dall'Agenzia delle entrate sulla fattura elettronica dovr\'a essere indicato:
\begin{itemize}
    \item per le operazioni interne, nel campo natura: si deve indicare il codice "N2.2 non soggette - altri casi", Riportando la dicitura "Operazione senza applicazione dell'IVA ai sensi dell'art.1, comma 58, lo. 190/2014";
    \item nel caso di operazioni con non residenti (es.: prestazion di servizi rese ad un soggetto passivo IVA sia UE che EXTRAUE), l'operazione \'e fuori campo IVA ai sensi dell'art. 7-ter, D.P.R. n.633/1972. In tale ipotesi si dovrebbe indicare il codice: "N2.1 - operazione non soggetta ad IVA ai sensi degli articoli da 7 a 7 septies, DPR 633/72";
    \item nel campo regime fiscale: indicare il codice RF19;
    \item nel campo descizione deve essere indicato che le somme non sono soggette a ritenuta
    \item l'imposta di bollo viene assolta compilando la sezione "dati bollo", valorizzando a SI, il campo "bollo virtuale", evidenziando in tal modo il debito per imposta di bollo (il riaddebito \'e facoltativo);
    \item che l'operazione \'e effettuata ai sensi dell'art. 1, comma 58 e 59, della Legge di stabilit\'a 2015 (Legge n.190/2014 e successive modifiche).
\end{itemize}

Conservazione in formato elettronico: l'emissione in formato elettronico comporta anche la necessit\'a di conservare le fatture con modalit\'a elettroniche ai sensi dell'art. 39, D.P.R. n. 633/1972. In tal caso, i soggetti obbligati potranni avvalersi del servizio di conservazione gratuita messo a disposizione dall'Agenzia delle entrate. 

Termini di trasmissione della fattura elettronica: in applicazione della disciplina di carattere generale, la fattura elettronica, ai sensi dell'art. 21, D.P.R. n. 633/1972, andr\'a trasmessa allo SDI:

\begin{itemize}
    \item entro 12 giorni dall'effettuazione dell'operazione nel caso di fattura immediata. Il termine decorre dal pagamento, nell'ipotesi di prestazione di servizi e dalla consegna del bene, nel caso di cessioni di beni (es.: commercianti al minuto/ingrosso)
    \item entro il giorno 15 del mese successivo, a quello di effettuazione dell'operazion nel caso di fattura differita
\end{itemize}

Nel caso di prestazioni di servizi, il momento di effettuazione dell'operazione ai fini IVA \'e la data di pagamento del corrispettivo ovvero fatturazione, se essa \'e anteriore al pagamento.

Nel caso di cessioni periodiche, la fatturazione deve avvenire al momento della data di pagamento dei corrispettivi ovvero fatturazione, se essa \'e anteriore al pagamento.

\subsection{Esoneri IVA}
I contribuenti forfetari sono esonerati:
\begin{itemize}
    \item dalla tenuta di registri contabili ed IVA relative registrazioni
    \item dalla liquidazioni periodiche
    \item dai versamenti periodici IVA (mensile)
\end{itemize}

\subsection{Obblighi IVA}
\begin{itemize}
    \item conservazione e numerazione delle fatture ricevute e delle bollette doganali
    \item certificazione dei corrispettivi e conservazione dei relativi documenti
    \item emissione di fatture in formato elettronico (XML)
    \item integrazione delle fatture relative agli acquisti intracomunitari superiori a 10.000 euro e per le altre operazoini di cui risultano debitori d'imposta quali il reverse charge interno, ad esempio ???

    \item presentazione dei modelli INTRASTAT (nei casi previsti);
    \item nel caso di passaggio di regime, rettifica alla detrazione (se dovuta);
    \item iscrizione all'archivio VIES;
    \item adozione del registratore telematico
\end{itemize}

Emissione nota di credito.

Trattamento imposta di bollo.

\subsection{Adozione del registratore telematico}
\subsection{Medici di medicina generale e foglio di liquidazione}

\subsection{Esterometro}
A partire dal primo luglio 2022, per i contribuenti forfettari obbligati all'emissione della fattura elettronica, sussiste l'obbligo di invio telematico "singolo" all'Agenzia delle entrate delle fatture estere nel formato XML.

\subsection{Indetraibilit\'a dell'IVA}
In considerazione del fatto che i contribuenti in regime forfetario non addebitano al cliente l'IVA per rivalsa, non \'e ammessa in detrazinoe l'IVA assolta, dovuta o addebitata sugli acquisti nazionali, comunitari e sulle importazioni.

\subsection{Assolvimento dell'imposta di bollo}
A decorrere dal primo luglio 2022, i contribuenti in regime forfettario obbligati all'emissinoe della fattua elettronica dovranno assolvere elettronicamente anche lìimposta di bollo dovuta sulle fatture emesse (qualora siano di importo superiore a euro 77,47). Il versamento dell'imposta di bollo deve essere effettuato entro l'ultimo giorno del secondo mese successivo alla chiusure del trimestre. 

La regola appena citata subisce una deroga per il secondo trimestre (aprile-giugno) dato che il versamento dell'imposta di bollo pu\'o essere effettuato entro il 30 settembre. Tali scandenze si applicano anche qualore l'imposta di bollo sia dovuta in misura non superiore a euro 250. Ne deriva che:

\subsubsection{Prospetto riassuntivo versamento imposta di bollo}
\begin{itemize}
    \item Bollo primo trim. superiore a 250 euro: entro il 31 maggio;
    \item Bollo primo trim. inferiore a 250 euro: entro il 30 settembre;
    \item Bollo secondo trim.: entro il 30 settembre;
    \item Bollo primo e secondo trim. inferiore a 250 euro: entro il 30 novembre;
    \item Bollo terzo trim. di qualsiasi importo: entro il 30 novembre;
    \item Bollo quarto trim. di qualsiasi importo: entro il 28 febbraio dell'anno successivo.
\end{itemize}

\subsubsection{Modalit\'a di pagamento}
Il versamento dell'imposta di bollo pu\'o avvenire o mediante il servizio presente nell'area riservata del sito dell'Agenzia delle entrate, con addebito su conto corrente bancario / postale o utilizzando il modello F24 utilizzando i sotto riportati codici tributo:

\begin{itemize}
    \item "2521" primo trimestre
    \item "2522" secondo trimestre
    \item "2523" terzo trimestre
    \item "2524" quarto trimestre
    \item "2525" SANZIONI
    \item "2526" INTERESSI
\end{itemize}

\subsubsection{Regime premiale e fatturazione elettronica}
Per la lLegge di stabilit\'a 2020, \'e previsto che per il contribuente in regime forfetario che per l'intero periodo d'imposta, emetta esclusivamente fatture elettroniche \'e riconosciuto il beneficio della riduzione di un anno del termine quinquennale di accertamento. La finalit\'a della norma \'e quella di incentivare l'utiolizzo della fattura elettronica, in luogo della fattura cartacea.

\subsubsection{Fattura elettronica PA}
Nei rapporti con la PA, per i contribuenti forfetari sussiste l'obbligo di emissione della fattura in formato elettronica. In tale ipotesi, si dovranno rispettare tutte le regole previste per la fattuarazione elettronica ed indicare sulla fattore il codice IPA della Pubblica amministrazione destinataria della fattura elettronica. Detto codice, di 6 caratteri, \'e l'unico indirizzo telematico a cui recapitare alle Pubbliche amministrazioni la fattura elettronica. 

\subsubsection{Delega alla fatturazione elettronica}
L'agenzia delle entrate ha aggiornato le speicigiche tecniche per l'attivazione delle deleghe agli intermediari ai fini della gestinoe dei servizi collegati alla fatturazione elettronica. Nel citato provvedimento \'e stato specificato che se, il soggetto delegante adotta il regime forfettario, gli elementi di riscontro sono l'importo del reddito lordo complessivo e l'importo corrsipondente al reddito assoggettato ad imposta sostitutiva nel quadro LM e l'importo corrispondente al reddito complessivo.

\subsection{Testo di riferimento}
Il testo di referimento per qualsiasi dubbio o problema \'e il "Testo Unico IVA".

\subsection{Come emettere fattura}
I dati necessari per una corretta emissione della fattura sono: partita iva destinatario, pec destinatario, codice fiscale, email, 
dettagli del prodotto o servizi, parametri fiscali (profilo fiscale) condizioni e informaizoni di pagamento, ed eventuali coordinate bancarie alla quale inviare il saldo.


\subsection{Quando emettere fattura}
Le principali situazioni in cui emettere fattura sono la vendita di beni e/o la prestazione di servizi. 

Per quanto riguarda la vendita di bene, la fattura deve essere emessa entro 12 giorni dalla data di effettuazione dell'operazione (quindi alla consegna o alla spedizione del bene), trattandosi quindi di fattura immediata,
oppure entro il 15 del mese successivo, se si fa riferimento a documenti di trasporto (DDT) o altri documenti che provano la consegna.

Nel caso di prestazioni di servizio, se la fattura \'e immediata, questa va emessa entro 12 giorni dalla data in cui \'e stata effettuata la prestazione.
Se il pagamento avviene in anticipo, la fattura va emessa alla data del pagamento anche solo parziale. 

Per i soggetti obbligati (la maggior parte delle partite IVA), la fattura deve essere inviata al Sistema di Interscambio (SdI) entro i termini sopraindicati. La data di emissione coincide con la data di trasmissione allo SdI.

Se si dovesse emettere una fattura in ritardo, si rischiano sanzioni amministrative che possono andare dal 90\% al 180\% dell'IVA non documentata correttamente. 

\subsection{Sistema di Interscambio}
Per Sistema di Interscambio si intende il canale telematico dell'Agenzia delle Entrate per la gestione della fatturazione eletteonica.

Il Sistema di Interscambio riceve le fatture elettroniche, effettua i controlli su questa e le invia ai destinatari. 

Il Sistema di Interscambio verifica che siano presenti le informazioni minime obbligatorie previste per legge 
(estremi identificativi del fornitore e del cliente, il numero, la data della fattura, la descrizione della natura, la quantit\'a, la qualit\'a del bene ceduto o del servizio perstato, l'imponibile, l'aliquota e l'IVA),
che i valori della partita IVA del fornitore siano esistenti, e quindi presenti in Anagrafe Tributaria, che sia compilato il campo "Codice Destinatario" e che ci sia coerenza tra i valori dell'imponibile, dell'aliquota e dell'IVA.

Se fosse presente un qualche tipo di errore, il SdI invia una notifica di scarto e la fattura viene considerata non emessa. 

Se la fattura \'e considerata corretta allora viene trasmessa all'indirizzo riportato nel campo "Codice Destinatario" e "PEC Destinatario" del file.

\subsection{Fatturazione: Guida AE}
Dal primo gennaio 2019 tutte le fatture emesse, a seguito di cessioni di beni e prestazioni di servizi effettuate tra soggetti residenti o stabilit in Italia, potranno essere solo fatture elettroniche. 

La fatturazione elettronica \'e obbligatoria sia nel caso di operazioni tra operatori IVA, sia nel caso di operazioni verso consumatori.

Le regole per predisporre, trasmettere, ricevere e conservare le fatture elettroniche sono definite nel provvedimento n. 89757 del 30 aprile 2018 pubblicato sul sito internet dell'Agenzia delle Entrate. 

Esistono regole diverse per le fatture emesse tra privati e fatture emesse verso le Pubbliche Amministrazioni. 

\subsection{Come predisporre, inviare e ricevere l'e-fattura}
Per compilare una fattura elettronica \'e necessario disporre di:
\begin{itemize}
    \item Un PC ovvero di un tablet o uno smartphone
    \item Un programma (software) che consenta la compilazione del file della fattura nel formato XML previsto dal provvedimento dell'Agenzia delle Entrate del 30 aprile 2018.
\end{itemize}

L'Agenzia delle Entrate mette a disposizione gratuitamente 3 tipi di programmi epr rpedisporre le fatture elettroniche:

\begin{itemize}
    \item una procedura web
    \item un software scaricabile su PC
    \item un'App per tablet e smartphone, denominata Fatturae
\end{itemize}

\subsubsection{Predisposizione della fattura elettronica mediante la procedura web dell'agenzia}
La prima volta che si utilizza la procedura occorre verificare i dati del fornitore (cedente/prestatore) che la procedura riporta in automatico recuperandoli dall'Anagrafe Tribuaria: il campo della Partita Iva non \'e modificabile, gli altri dati possono essere variati. 

Una volta salvati i dati riportati nella schermata, la stessa non verr\'a pi\'u riproposta per la compilazione delle successive fatture. 

Quindi, inserire i dati del clienti (cessionario/committente), ricordandosi di compilare sempre il campo "Codice Destinatario": quest'ultimo campo potr\'a essere compilato con il codice di 7 cifre alfanumerico che avr\'a comunicato il cliente e rappresenta l'indirizzo telematico dove recapitare le fatture. 

Se il cliente dovesse comunicare un indirizzo PEC (quale indirizzo telematico dove intende ricevere la fattura), il campo "Codice Destinatario" dovr\'a essere compilato con il valore "0000000" e, nel campo "PEC destinatario", andr\'a riportato l'indirizzo PEC comunicato dal cliente. 

Se il cliente non comunica alcun indirizzo telematico, sar\'a sufficiente compilare solo il campo "Codice Destinatario" con il valore "0000000".

Attenzione: se il fornitore inserisce solo il valore "0000000" nel campo "Codice Destinarario", il SdI non riuscir\'a a consegnare la fattura elettronica al cliente, ma la metter\'a a disposizione di quest'ultimo in un'apposita area di consultazione riservata del sito dell'Agenzia. Quindi, sar\'a importanto che il fornitore consegni al cliente una copia, anche su carta, ricordandogli che la fattura originale \'e quella elettronica e che potr\'a consultarla e scaricarla dalla sua area riservata del sito dell'Agenzia delle Entrate. 

Inserire i dati relativi alla natura, quantit\'a e qualit\'a del bene ceduto o del servizio prestato, nonch\'e i valori dell'imponibile, dell'aliquota Iva e dell'imposta (ovvero, nel caso di operazioni esenti, non imponibili ecc., l'apposito codice che identifica la "natura" dell'operazione ai fini IVA).

\subsubsection{Invio fattura elettronica al cliente}
Le fatture elettroniche vanno sempre inviate ai propri clienti attraverso il SdI (Sistema di Interscambio), altrimenti sono considerate non emesse. 

Per trasmettere al SdI il file XML della fattura elettronica ci sono diverse modalit\'a:

si pu\'o utilizzare un servizio online presente nel portale "Fatture e Corrispettivi" che consente l'upload del file XML preventivamente predisposto e salvato sul proprio PC;

si pu\'o utilizzare la procedura web ovvero l'App Fatturae messe a disposizione gratuitamente dall'Agenzia delle Entrate;

si pu\'o utilizzare una PEC, inviando il file della fattura come allegato del messaggio di PEC all'indirizzo "sdi01@pec.fatturapa.it"

si pu\'o utilizzare un canale telematico (FTP o Web Service) preventivamente attivato con il SdI.

\subsubsection{Cosa fa il Sistema di Interscambio quando riceve una fattura}
\begin{itemize}
    \item verifica che siano presenti almeno le informazioni minime obbligatorie previste per legge
    \item verifica che i valori della partita Iva del fornitore (cedente/prestatore) e della partita Iva oppure del Codice Fiscale del cliente (cessionario/committente) siano esistenti, cio\'e presenti in Anagrafe Tributaria
    \item verifica che sia inserito in fattura l'indirizzo telematico dove recapitare il file
    \item verifica che ci sia coerenza tra i valore dell'imponibilem dell'aliquota e dell'Iva
\end{itemize}

\subsubsection{Reverse charge}
Il reverse charge, o inversione contabile, \'e un particolare meccanismo di applicazione dell'imposta sul valore agginuto, per effetto del quale il destinatario di una fornitura di beni o prestazione di servizi (cliente), se soggetto passivo nel territorio dello Stato, \'e tenuto all'assolvimento dell'imposta in luogo del fornitore o prestatore.

L'inversione contabile \'e uno strumento utile per l'erario e tale meccanismo prevede che il cedente/prestatore riceva dal cliente esclusivamente l'importo imponibile del bene ceduto o della prestazione eseguita, con la conseguenza che non si configura in capo ad esso l'obbligo di versare l'Iva dell'operazione eseguita. 

Il settore maggiormente interessato al meccanismo dell'inversione contabile \'e quello edile, soprattutto per ci\'o che concerne i rapporti tra subappaltatori o tra subappaltatore e appaltatore.

\subsubsection{Fattura classica}

\subsubsection{Fattura semplificata}
Data di emissione del documento e il numero progressivo per identificarla in modo univoco, i dati del prestatore (numero di partita Iva, denominazione o raginoe sociale, dati anagrafici(nome, cognome, residenza)), Codice fiscale opartita Iva del cliente, se stabilito in altro Stato dell'Unione europea servir\'a il numero di identificazione Iva di quel dato Paese, descrizione breve dei beni o dei servizi ceduti, l'ammontare totale del costo dell'operazione e l'aliquota Ive che viene aggiunta. 

\subsubsection{Software per fatturazione}
\begin{itemize}
    \item Fattura24
    \item Aruba Fatturazinoe Elettronica
    \item Fatture in Cloud
    \item Libero SiFattura
    \item Invoice Buddy
    \item Fattura per tutti
    \item Fiscozen (Gestione fiscale con supporto consulenziale)
    \item FatturaElettronica APP
    \item Danea EasyFatt
    \item Agenzia delle Entrate
    \item AssoInvoice
    \item SDiPEC
\end{itemize}

\section{Contratti}

il software viene considerato opera creativa

diritto morale si distingue in diritto morale e diritto patrimoniale

per copyright ci si rifericsce al diritto patrimoniale. In inglese si parla anche di royalties

il diritto d'autore nasce con la creazione dell'opera
pu\'o coinvolgere pi\'u autori
\'e consigliabile depositare l'opera presso un apposito ente (SIAE)
l'opera verr\'a contrassegnata con la c cerchiata

Fonti normative
Legge 633/41 deve essere integrata
articoli 2575-2584 del Codice Civile
Legge 248/2000
Decreto Legge 72/2004
Regolamento in materia di tutela del diritto d'autore sulle reti di comunicazinoe elettronica

Cosa garantiscono le normative aggiuntive?

Internet ha portato a qualcosa di non tangibile, di immateriale

\'E reato condividere tramite programmi di file sharing opere coperte da copyright

\'E possibile ordinare la rimozioni di contenuti pubblicati in rete o, nei casi pi\'u gravi la rimozione di intere piattaforme
(Regolamente in materia di tutela del diritto d'autore sulle reti di comunicazinoe elettronica)

Tutela del d.D'autore
sono previste
azioni penali
azioni civili

per commettere reato
intenzionalit\'a
profitto (vantaggio economico o vantaggio morale)
Profitto economico: maggior guadagno o minor spesa

\subsection{Introduzione}
Il diritto privato si divide in diritto civile, diritto commerciale, diritto del lavoro, diritto agrario.

Una branca del diritto civile \'e quella dei diritti di propriet\'a intellettuale a cui fanno capo tre grandi aree: diritto d'autore, diritto dei brevetti e diritto dei marchi.

Il diritto d'autore \'e disciplinato principalmente dalla Legge n. 633 del 1941, con successive modifiche e protegge le espressioni artistiche, nonch\'e beni immateriali creativi di carattere utilitaristico come software e banche dati.

Il diritto d'autore protegge opere dell'ingegno creativo ed originali in campo letterario, artistico, musicale, cinematografico, software, ecc. Questo nasce automaticamente al momento della creazione dell'opere, anche se \'e consigliata la registrazione presso enti come la (SIAE). In generale vale fino a 70 anni dopo la morte dell'autore e pu\'o proteggere opere come un romanzo, una canzone, un dipinto, un film o un codice software originale.

La propriet\'a industriale protegge invenzioni, segni distintivi, disegni industriali e quindi elementi usati per attivit\'a economiche/industriali. La propriet\'a industriale nasce solo dopo la registrazione presso enti competenti che possono rilasciare ad esempio brevetti, marchi o design. I brevetti durano per massimo 20 anni, i marchi sono rinnovabili ogni 10 anni e il design pu\'o durare al massimo 25 anni. Esempi di opere protette da propriet\'a industriale sono il logo Nike (marchio), la forma di un bottiglia particolare (design), una nuova tecnologia per motori (brevetti).

Nel mondo del digitale, il diritto d'autore ha assunto un ruolo centrale data la facilit\'a di copia e diffusione delle opere su internet. I software vengono protetti come opere lettarare e gli strumenti di tutela del diritto d'autore sono i watermark, i DRM, e le licenze Creative Commons.

\subsection{Classificazione dei diritti}
Esistono i diritti patrimoniali ed i diritti non patrimoniali. I primi sono quelli che hanno ad oggetto un bene suscettibile di valutazione economica, come ad esempio un immobile. I secondi sono quelli non soggetti a valutazione economica, come ad esempio i diritti della personalit\'a, il diritto alla famiglia, il diritto alla vita ecc.

Il diritto assoluto, fatto valere nei confronti di qualsiasi persona. 

I diritti relativi, fatti valere soltatno nei cofronti di una persona. 

I diritti reali: hanno ad oggetto le cose, in maniera diretta, non mediato. Questi sono assoluti (erga omnes), il diritto di propriet\'a ad ese.

Diritto di obbligazione: esigere una prestazione da una determinata persona.

Diritti potestativi: sono poteri diretti a creare, modificare, estinguere una sitaizone mediante un atto unilaterale. Chi la riceve, la pu\'o soltanto accettare (il diritto di ritirarsi da un contratto).

\subsubsection{Classificazione dei diritti d'autore}
Il diritto d'autore si articola in diritti morali e diritti di utilizzazione economica dell'opera. 

I primi sono inalienabili, irrinunciabili e imprescrittibili, nonch\'e indipendenti dai diritti patrimoniali. Essi infatti attengono alla sfera dei c.d. diritti della personalit\'a. Sono il diritto a rivendicare la paternit\'a dell'opera, il diritto di opporsi a ogni deformazinoe e mutilazione e/o modificazione che possa apportare un pregiudizoi all'onore e alla reputazione dell'autore e il diritto a ritirare l'opera dal commercio quando ricorrano "gravi ragioni morali"; il diritto di inedito. Questi diritti possono essere esercitati anche dagli eredi dell'autore.

I secondi invece sono rinunciabili e cedibili a terzi, e hanno una durata stabilita dalla legge e si estinguono infatti 70 anni dopo la morte dell'autore; sono inoltre tra loro indipendenti. Sono il diritto di riproduzione in pi\'u esemplari dell'opera; il diritto di trascizione dell'opera orale; il diritto di esecuzione, rappresentazione o recitazione in pubblico; il diritto di comunicazione; il diritto di distribuzione; il diritto di elaborazione, di traduzione e di pubblicazione delle opere in raccolta; il diritto di onleggio e di dare in prestito.

\subsection{Definizione di software sotto il profilo giuridico}
Secondo la definizione adottata dal WIPO (World Intellectual Property Organization) nel 1984, il software consiste “nell’espressione di un insieme organizzato e strutturato di istruzioni (o simboli) contenuti in qualsiasi forma o supporto (nastro, disco, film, circuito), capace direttamente o indirettamente, di far eseguire o far ottenere una funzione, un compito od un risultato particolare per mezzo di un sistema di elaborazione elettronica dell’informazione”.

Si distingue tra programmi di base e programmi applicativi. I primi sono la base sulla quale \'e possibile eseguire le operazioni pi\'u comuni, la seconda sulla base della quale sono idonei ad esercitare una funzione per l'utente finale.

Il software \'e costituito da codice sorgente e codice oggetto: il primo \'e il linguaggio convenzionale del programmatore per scrivere il software, il secondo \'e il linguaggio della macchina, ottenuto per traduzione del primo, necessario per essere eseguito dalla macchina stessa. L'operazione inversa, cio\'e di ri-traduzione del codice oggetto in codice sorgente, \'e una complicata operazione di reverse engineering. 

Il codice sorgente rimane segreto o non disponibile nelle consuete licenze d'uso di software. 

Il software appartiene alla categoria dei beni giuridici immateriali, non soggetto ad usura e che pu\'o essere utilizzato da un numero indefinito di soggetti, senza che la sua utili\'a venga diminuita. Pertanto il software rientra nelle creazioni intellettuali ed in quanto tale \'e tutelabile come opera d'ingegno, soggetta a diritto d'autore e come invenzione industriale, soggetta a brevetto. (art. 2575 c.c. e art 2585 c.c., rispettivamente). 

\subsection{Diritto d'autore (copyright)}
Il software essendo assimilato ad un'opera letteraria, rientra tra le opere protette ai sensi della Legge sul diritto d'autore. 

La tutela apprestata dalla LdA implica che venga protetta la forma espressiva del software, non le idee (o funzioni) sui cui il programma \'e basato. Pertanto il diritto d'autore tutela il codice sorgente, il codice oggetto, il materiale preparatorio, ma non le idee ed i principi alla base del codice sorgente od oggetto di un programma.

Il diritto d'autore sul software si dividie in una componente morale, che rimane sempre attribuita all'autore del programma, e in una di natura patrimoniale, che pu\'o essere oggetto di cessione. Entrambe nascono in modo originario al momento della creazione del programma come previsto dall'art. 6 LdA, senza alcuna formalit\'a. 

I diritti morali sono inalienabili ed esistono indipendemente dalla cessione a terzi dell'opera o dalla stipula di un contratto che ne preveda l'utilizzazione economica in quanto sono strettamente legati all'autore. Tali diritti, sanciti dagli artt. 20 e 24 della LdA sono: il diritto alla paternit\'a del software, il diritto all'onere ed alla reputazione, il diritto di inedito, che concede all'autore la possibilit\'a di scegliere se pubblicare o meno il software.

I diritti patrimoniali permettono invece all'autore lo sfruttamente economico dell'opera (previsto dall'art. 12 LdA), sia in forma originale che nella forma derivata, e di ricevere un compenso per ogni tipo di utilizzo della stessa, tra cui: attivit\'a di commission/sviluppo software, il noleggio, la licenza d'uso.

Per quanto riguarda i diritti patrimoniali, il diritto d'autore sul software si differenzia sostanzialmente da quello classico. Sono soggette all'autorizzazione del titolare dei diritti anche il caricamento, la visualizzazione, l'esecuzione, la trasmissione o la memorizzazione se comportano una riproduzione del programma, nonch\'e la traduzione, adattamento, trasformazione e ogni altra modificazione del programma per elaboratore e la riproduzione dell'opera che risulti.

Il diritto permette comunque l'osservazione solo dall'esterno del funzionamento del programma (Black box analysis).

Nel campo della tutela dei programmi per elaboratore, la decompilazione pu\'o essere effettuata solo nel rispetto di precisi limiti:
\begin{itemize}
    \item la decompilazione deve essere eseguita da chiha il diritto di usare il software
    \item le informazioni non devono essere gi\'a facilmente reperibiki e rapidamente accessibili
    \item la decompilazione deve essere limitata solo alle parti del programma interessate per conseguire l'interoperabilit\'a
    \item le informazioni acquisite nonn devono essere comunicate a terzi
    \item le informazioni acquisite non possono essere utilizzate per costruire programmi sostanzialmente simili nella loro forma espressiva
\end{itemize}

\subsection{Brevettabilit\'a del software}
L'invenzione industriale \'e tutelata attraverso il brevetto. L'art. 12 della Legge sulle Invenzioni (L.I) dispone che i programmi di elaboratori "in quanto tali" non sono considerati come invenzioni (e pertanto non sono brevettabili).

Il fatto che non si possa brevettare il codice in quanto tale, non esclude il brevetto di un'invenzione di combinazione, in cui si faccia uso di un elaboratore programmato in modo da ottenere un risultato tecnico nuovo e originale tramite il coordinamento di elementi e mezzi gi\'a conosciuti.

Il software deve necessariamente rispettare i requisiti richiesti per qualsiasi invenzione, ovvero:
\begin{itemize}
    \item carattere tecnico
    \item novit\'a
    \item originalit\'a
    \item liceit\'a
    \item applicabilit\'a industriale
\end{itemize}

Per brevettare un software \'e ncesasrio presentare domanda di brevetto presso l'Ufficio Centrale Brevetti (UCB). I diritti sull'invenzione durano venti anni, al decorrere dei quali, il brevetto non pu\'o pi\'u essere rinnovato n\'e esserne prorogata la durata.

\subsection{Diritto d'autore e brevetto a confronto}
Il brevetto ed il diritto d'autore (copyright) sono due strumenti notevolmente diversi per garantire la tutela giuridica del software, in quanto implicano differenti modi di acquisizione dei diritti e della durata degli stessi. 

Il brevetto protegge l'idea, non permette la riproduzione e vendita di un programma dotato di funzionalit\'a  analoghe o identiche. Il diritto d'autore protegge la forma, a prescindere del suo contenuto, pertanto se ne vieta la duplicazione del codice, ma non l'idea o la funzionalit\'a cui si riferisce quel codice.

Nel caso del brevetto, per\'o, il codice sorgente entra a far parte del patrimonio scientifico della collettivit\'a, diventa quindi pubblico, dal momento che per essere ottenuto, \'e necessario pubblicare il codice sorgente stesso. 

I requisiti per la registrazione di un brevetto sono industrialit\'a, novit\'a, originalit\'a e liceit\'a, mentre il diritto d'autore sorge automaticamente in capo all'autore al momento dell'estrinsecazione dell'idea nella realt\'a fattuale.

I diritti sull'invenzione, il brevetto quindi, laddove ne ricorrano i requisiti, sorge al momento della registrazione del brevetto stesso, mentre il diritto d'autore sorge alla data di creazione dell'opera stessa.

Il brevetto ha un termine di validit\'a di 20 anni dalla domanda di deposito, mentre il diritto d'autore ha una durata di 70 anni dalla morte dell'autore.

In ogni caso, bisogna sempre valutare gli obiettivi e le strategie del titolare del software da dover soddisfare a fronte di una scelta per la forma di tutela del software. 

\subsection{Distribuzione del software. Principio di esaurimento comunitario}
Il principio di esaurimento comunitario afferma che il titolare di Il titolare di un diritto di proprietà industriale (come un marchio, un brevetto, o un diritto d’autore) ha il controllo esclusivo sulla prima immissione sul mercato dei prodotti coperti da quel diritto.

Ma, una volta che il titolare mette in commercio (direttamente o con il suo consenso) un bene all’interno dell’UE, non può più opporsi alla sua successiva rivendita o distribuzione all’interno del mercato europeo. In pratica, il suo diritto si “esaurisce”.

Quindi, il cessionario del software ha il diritto di trasferire a sua volta a terzi il software, senza necessit\'a di autorizzazione da parte del suo titolare originario.

In particolare, come dal caso Oracle, il diritto di distribuzinoe della copia di un programma per elaboratore \'e esaurito qualori il titolare abbia conferito, dietro pagamento di un compenso, il diritto di utilizzare la copia del software senza limitazioni di durata, a fronte di un compenso corrispondente al valore economico della copia stessa.  In questo caso, infatti, il titolare dei diritti ha già soddisfatto la propria rendita di posizione monopolistica con la percezione del corrispettivo economico versato dal licenziatario per la licenza, e non ha quindi più diritto a controllare l’eventuale successiva commercializzazione dell’opera. Inoltre, si \'e ritenuto che l'esaurimento del diritto di distribuzione \'e efficace indipendentemente dal fatto che la vendita riguardi una copia tangibile o intangibile del programma, ovvero la vendita di un programma per elaboratore su CD-ROM o DVD e la vendita medianete download via Internet, quanto agli effetti, sono operazioni negoziali analoghe. 

Quindi, in caso di rivendita di una licenza di utilizzazione, ogni acquirente successivo della licenza pu\'o avvalersi dell'esaurimento del diritto di distribuzione e dunque essere considerato quale leggittimo acquirente di una copia di un programma per elaboratore. 

In conclusione, qualora un software sia concesso in "licenza" per un periodo di tempo illimitato e a fronte di un pagamento, \'e legittima la licenza di software usato (c.d. relicensing) da parte dell'acquirente stesso, sia nel caso di vendita di software tramite un supporto fisico, sia riptrodoot online. Ci\'o significa che le eventuali restrizioni del titolare del software dirette a impedire l'ulteriore circolazione del software, CONTENUTE NEL CONTRATTO DI LICENZA, sono invalide, in quanto predisposte al fine di aggirare l'esaurimento del diritto di distribuzione, ai sensi dell'art. 1379 c.c.

Tuttavia, il software acquistato attraverso un contratto on-line, non avendo supporto fisico che lo racchiude, rimane inscindibilmente legato all'elaboratore elettronico sul quale \'e installato; quindi, nonostante con l'acquisto del programma si sia verificato l'effetto proprio dell'esaurimento del diritto, l'acquirente, in virt\'u dei limiti imposti dalla normativa in tema di programma per elaboratore, non pu\'o effettuare una copia del programma, diversa da quella di back-up. 

Ne deriva che, una volta perfezionato l'acquisto del software ed installato lo stesso sul proprio computer, l'acquirente, per commercializzarlo nuovamente, deve necessariamente rendere inutilizzabile la propria copia (originaria) al momento della rivendita, rendendo inutilizzabile sul proprio elaboratore, in modo definitivo, irreversibile e verficiabile la copia origniariamente acquistata. Solo in tal modo infatti viene evitata l'immissione in commercio di un numero di esemplari del programma superiore a quelli originariamente ceduti dal titolare. 

Il titolare del dirtito d'autore potr\'a verificare che effettivamente l'acquirente iniziale abbia smesso di utilizzare il programma dopo averlo venduto utilizzando delle misure tecniche di protezione, quali le chiavi di accesso al software, che agiscono on line direttamente sul singolo elaboratore connesso in rete.

Connesso all'obbligo del rivenditore di rendere inutilizzabile il programma, poi, \'e il divieto di scorporare le licenze. Ovvero \'e stato stabilito il principio di inscinbilit\'a della licenza del software, ovvero, qualora la licenza acquisita dell'iniziale acuqirente prevedea un numero di utenti superiore alle sue esigenze, questi, per effetto della regola dell'esaurimento.

Infine l'acquire di un software usato non pu\'o beneficiare dell'assistenza e della manutenzione allo stesso, in quanto, come chiarito dalla Corte di Giustizia, l'erogazione di tali servizi costituisce parte integrante della copia inizialmente scaricata dal primo acquirente. Per usufruire anche di questi servizi, il successivo acquirente dovr\'a espressamente richiederli. 

Nonostante tutte le problematiche per le aziende produttrici di software determinate dal principio di esaurimento, mantenere il controllo dei propri software dopo la prima vendita \'e possibile, regolamentendo con attenzione la cessione dei diritti nel contratto di contratto di licenza d'uso del software o di sviluppo del software. A tale scopo, \'e quindi opportuno disciplinare dettagliatamente il contenuto dei diritti ceduti, LIMITANDO NEL TEMPO IL DIRITTO D'USO, prevendendo NON un corrispettivo unitario, bens\'i un canone periodico, che non sia equivalente al valore economico di mercato della copia del software, cos\'i da assimilare il relativo contratto di licenza, non ad una vendita, ma ad una LOCAZIONE. 

Un'altra possibilit\'a, sempre pi\'u utilizzata dai vari produttori di software, \'e il cloud computing, nella quale non \'e necessario cedere una copia del codice sorgente ne un a copia del codice oggetto, e che il contratto di licenza \'e in tal modo assimilabile ad un abbonamento piuttosto che ad un contratto di locazione. 

\subsection{Contratto di sviluppo del software}
Il contratto di sviluppo del software \'e un contratto atipico (cio\'e non disciplinato dalla legge) con cui un soggetto (committente) richiede la realizzazione di un determinato software personalizzato (custom made), cio\'e un software avente specifiche caratteristiche, richiesta dal committente stesso, ad un altro soggetto, professionista autonomo o lavoratore dipendente (sviluppatore), verso un determinato corrispettivo. 

Il contratto di sviluppo del software pu\'o essere strutturato come un contratto d'opera (regolato dall'art.2222 e ss.c.c.) o come un contratto di appalto (regolato dall'art, 1655 c.c.). Tale diversa configurazione si riflette direttamente sul regime di responsiablit\'a dello sviluppatore e sui diritti di sfruttamente del software. 

Il contratto di sviluppo software ricondotto come contratto d'opera rende lo sviluppatore responsabile solo per danni derivanti da dolo o colpa grave, in base all'art. 2236 c.c.

Il contratto di sviluppo software ricondotto come contratto d'appalto rende lo sviluppatore obbligato ad un risultato, ed \'e quindi tenuto a consegnare al committente il software secondo specifiche richieste ed esigenze precisate da quest'ultimo. In tal caso, trova applicazione la garanzia per vizi e difformit\'a di cui agli artt. 1667 e 1668 c.c.

I diritti morali rimangono allo sviluppatore, in quanto autore del bene e poich\'e inalienabili; dunque egli rimane titolare dell'opera e ne mantiene la possibilit\'a di rivendicare la paternit\'a dell'opera in qualsiasi momento e di vietarne eventuali modifiche. 

I diritti patrimoniali sul software spettano anch'essi allo sviluppatore, il quale potr\'a modificare, riprodurre, distribuire il software a terze o cederlo ad altri committenti. Il committente godr\'a di una semplice licenza d'uso del software. 

Nel caso, per\'o, di rapporto di lavoro subordinato, tra committente e sviluppatore, i diritti di utilizzazione economica ricadono sul datore di lavoro.

La conclusione di un contratto di sviluppo software si articola generalmente in tre fasi: studio di fattibilit\'a, progettazione e collaudo.

Di centrale importanza \'e il piano di fattibilit\'a iniziale, nel quale entrambe le parti si impegnano a soddisfare le necessit\'a dell'uno e dell'altro, cercando di contenere al massimo la possibilit\'a di contenziosi  in fase di esecuzione del contratto. In questa fase, \'e gi\'a importante l'apporto dello sviluppatore ed \'e quindi fondamentale la cooperazione da ambedue le parti.

Altrettanto importante \'e la fase di collaudo, all'interno della quale sar\'a necessario precisare quali test, in che tempi ed in che forme dovranno essere effettuati, quali i casi di fallimento dei test stessi ed in che termini dovr\'a avvenire la comunicazione del mancato superamento di uno o pi\'u test. Il tutto dovr\'a essere descritto in una apposita procedura di accettazione del software, nel quale vengono riportate anche le modalit\'a alla quale sottoporre i test e quali i comportamenti attesi. 

In taluni casi, sar\'a necessario fare anche riferimento al codice sorgente stesso, se necessario, da entrambe le parti, sia sviluppatore che committente.

In generale, un contratto di sviluppo del software deve rispettare i seguenti aspetti: 
\begin{itemize}
    \item obblighi del committente, in relazione alle richieste ed agli obiettivi del software
    \item fasi della realizzazione del software
    \item obblighi dello sviluppatore
    \item modifiche al software in corso d'opera e successive alla consegna
    \item termini e tempisitche della denuncia di eventuali vizi e malfunzionamenti del software
    \item eventuale esclusiva di utilizzo del software da parte del committente
    \item propriet\'a intellettuale e del diritto d'autore sul software
    \item durata della garanzia sul software
\end{itemize}

\subsubsection{Titolarit\'a del codice sorgente}
Esistono diverse ipotesi circa la titolarit\'a del codice sorgente nel contratto di sviluppo software. 

\subsubsection*{Ipotesi A}
Il Fornitore cede sia il codice sorgente dell'applicazione cos\'i come la relativa documentazione tecnica.

Il Committente conseguie il diritto di modificare ed estendere il software secondo le proprie esigenze; inoltre il committente acquisisce ogni diritto patrimoniale sullo stesso, o su una specifica parte di questo.

Il Fornitore si impegna altres\'i a risarcire e tenere indenne il Committente da qualsivoglia azione che dovesse essere intrapresa da terzi in relazione a presunti diritti vantati sul software; nonch\'e a intervenire nei giudizi civili e/o penali eventualmente promossi da terzi, anticipando spese e oneri che il Committente si trovasse a dover affrontare in relazione a detti giudizi.
le
\subsubsection*{Ipotesi B}
Il Fornitore cede al Committente il codice oggetto, il codice sorgente e la relativa documentazione tecnica. 

Il Committente consegue il diritto di modificare ed estendere il software secondo le proprie esigeneze, ma si impegna a NON CEDERE A TERZI IL CODICE SORGENTE E LA DOCUMENTAZIONE TECNICA. Inoltre, il Committente \'e obbligato a destinare il software consegnatogli dal Fornitore e le sue eventuali modifiche ed estensioni successive ad un mero uso interno.

Il Fornitore conserva i diritti patrimoniali e morali del software sviluppato.

\subsubsection{Ipotesi C}
Il Fornitore trasferisce al Committente il solo diritto di utilizzo del software in versione oggetto. 

Il codice sorgente del software rimane nella esclusiva disponibilit\'a del Fornitore. Quest'ultimo conserva in capo a s\'e la titolarit\'a di ogni diritto patrimoniale.

\subsubsection{Ipotesi C1}
Il Fornitore trasferisce al Committente il solo diritto di utilizzo del software in versinoe oggetto.

Il codice sorgente del software rimane nella esclusiva disponibilit\'a del Fornitore. Quest'utlimo conserva in capo a s\'e la titolarit\'a di ogni diritto connesso allo sfruttamenteo commerciale del softwaer sviluppato; tuttavia il Fornitore si impegna a non cedere a terzi il software, sia sotto forma di sorgente, sia nella versione eseguibile, per .... mesi dalla sua consegna al committente.

\subsubsection{Ipotesi D (Code Escrow)}
Il Fornitore trasferisce al cliente il solo diritto di utilizzo del programma in versione oggetto. 

Il Fornitore rimane titolare esclusivo del codice sorgente, ma si impegna a proprie spese, a depositarlo,, unitamente alla relativa documentazinoe tecnica presso ...., affinch\'e quest'ultimo lo custodisca e lo consegni al Committente nelle ipotesi di fallimento o sottoposizione del Fornitore ad altra procedura concorsuale, nonch\'e di cessazione dell'attivit\'a del medesimo. 

Il terzo depositario si impegna a conservare il codice sorgente e la rielativa documentazione tecnica per .... anni; al termine di questo periodo, egli provveder\'a alla distruzione materiale del codice stesso e della documentazione tecnica.

\subsection{Licenza d'uso software}
In linea generale, e salvo diversa pattuizione tra le parti, il diritto di propriet\'a e i relativi diritti di utilizzazione economica in ordine sia al codice sorgente che al codice oggetto spettano al soggetto committente, in base alle norme in tema di contratto d'opera.

Tuttavia \'e frequente che il fornitore ceda al committente solo il programma eseguibile, quindi solo il codice oggetto, ma non il codice sorgente, che consentirebbe di modificarne la struttura ed il suo funzionamento. In questo caso, il committente \'e soltanto licenziatario del codicec sorgente; quest'ultimo resta di esclusiva propriet\'a dello sviluppatore  software che ne detiene tutti i relativi diritti morali e patrimoniali, mentre al committente pu\'o essere vietato di effettuare alcuna forma di cessione, distribuzione, diffusione, adattamento, trasformazione ed ogni altra modifica o intervento al codice sorgente, oppure essere concesso di utilizzarlo internamente per eventuali modifiche ed estensioni. 

Questo pu\'o essere problematico quando la software house non \'e pi\'u in grado di fornire manutenzione ed assistenza del software, perch\'e venuta meno o per qualsiasi altra ragione, non consentendo cos\'i al committente-cliente di effettuare aggiornamenti, migliorie o nuovi interventi sul software. 

Per risolvere questo problema, si pu\'o ricorrere al c.d. contratto di source code escrow, con la quale la software house si occupa di rilasciare una copia ad una terza parte di fiducia, anche verso il cliente-committente stessa, cos\'i nel caso di impossibilit\'a di adempienze da parte della software house, il committente pu\'o richiedere copia del codice all'escrow.

Ritornando al contratto di licenza d'uso del software, questo \'e un tipo di contratto che trasferisce al cliente tutti o parte dei propri diritti patrimoniali e d'autore, per corresponsione di un determiato canone d'uso, periodico o annuale.

I diritti patrimoniali che nell'ambito di un contratto di licenza d'uso, l'autore del programma pu\'o trasferire al beneficiario sono elencati agli artt. 12 e ss. LdA, ovvero:

\begin{itemize}
    \item diritto di pubblicazione
    \item diritto di riproduzione
    \item diritto di modificazione, trasformazione e traduzione dell'opera
    \item diritto di distribuzione
    \item diritto di noleggio e prestito
\end{itemize}

Quindi, il licenziante, rimane comunque proprietario del software stesso, non come in un contratto di cessione del software in cui il cliente ne ottiene tutti i diritti patrimoniali, ivi compreso il codice sorgente. 

In questo caso, quindi, il codice oggetto del programma, verr\'a fornito al cliente committente, ma non il suo codice sorgente. Il codice oggetto in questione potr\'a essere soltanto eseguito o utilizzato, ma non modificato. In tal senso il software viene definito "proprietario", concedendo quindi soltanto la facolt\'a di installazione e utilizzo.

\subsection{Trial License}
Una licenza che funge da prova per il licenziatario, il quale pu\'o scaricare il software senza dover apporre alcuna firma e utilizzarlo per un determinato periodo di tempo, alla scadenza del quale valuter\'a se acquistarlo o meno.

Il contratto di licenza d'uso del software deve indicare chiaramente l'ambito di utilizzo del software da parte del licenziatario, cio\'e le modalit\'a, i contesti e i limiti entro i quali il licenziatario \'e autorizzato a utilizzare il software in modo da tutelare i diritti del licenziante ed evitare abusi o utilizzi impropri del software.

Occorre quindi definire:
\begin{itemize}
    \item Il numero di utenti che possono usufruire della licenza d'uso di software
    \item I tipi di dispositivi
    \item L'utilizzo del software, che pu\'o essere commerciale e/o personale
    \item Ambito geografico in cui pu\'o essere utilizzato il software
    \item se il licenziatario \'e autorizzato o meno a modificare o personalizzare il software ed in quali limiti
    \item quali sono le limitazioni e i divieti d'uso del software in capo al licenziatario, come ad es. il divieto di reverse engineering, di distribuire il software a terze parti, di concedere sub-licenze, di utilizzare il software per scopi illegali o contrari all'etica.
\end{itemize}

In caso di cessione di ramo d'azienda, la giurisprudenza prevalente ritiene che l'acquirente dell'azienda non subentri automaticamente nel contratto di licenza del software ai sensi dell'art. 2558 c.c., ma solo se ci\'o sia espressamente previsto nel contratto di cessione, dato il carattere personale dei diritti sul software, legato all'inventiva, che implica il pieno dominio del titolare anche nella determinazione dell'uso da parte dei terzi, salvo una diversa pattuizione tra le parti.

\subsection{Licenze d'uso Open Source}
Per venire incontro al sapere informatico, si \'e sviluppato un secondo modello, denominato open source (OS), anch'esso ancorato al diritto di autore, il quale prevede la circolazione/distribuzione del software nella sua versione in codice sorgente. 

Al licenziatario, inoltre, vengono espressamente attribuite facolt\'a e diritti ulteriori rispetto alla mera utilizzazione della singola copia del softwaer concessa in licenza; in particoalre, il licenziatario pu\'o alle condizioni poste dalla licenza d0uso riprodurre il programma, modificarlo, effettuarne copie e talvolta anche distribuirlo, il tutto senza corrispettivo. 

Affinch\'e si possa parlare di software OS, occorre che il codice sorgente del progrmama per elaboratore sia accessibile a chiunque, o in forma diretta, oppure deve essere ben pubblicizzato il modo per ottenerlo. 

L'essenza dell'OS consiste nella mancanza di qualunque restrizione circa l'uso del programma, che pu\'o essere anche ridistribuito dall'utente a terzi e ci\'e senza previsinoe di alcun corrispettivo. 

Sotto licenza OS, \'e consentito a tutti operare le modiche del programma in quanlunque forma e modo.

Le licenze OS prevedono il diritto non esclusivo, non soggetto a limitazioni temporali o territoriali, di utilizzare, distribuire, riprodurre, modificare, ridistribuire e rappresentare in pubblico il programma per elaboratore. Anche nel caso di sublicenza, tali facolt\'a sono attribuite al sub licenziatario, il cui successivo utilizzo viene regolato dalle condizioni e dai termini espressi dalla licenza originaria.

Ne caso in cui poi il software OS venga incluso in un altro software per dar vita ad un prodotto diverso da distribuire al pubblico, la parte di software OS dovr\'a essere sempre regolata dalla licenza originaria. Per le opere derivate, inoltre, dovr\'a essere predisposta una versione diversa dell'originario schema contrattuale della licenza, omettendo i riferimenti al soggetto che per prima ha realizzato e concesso in uso il programma o d altrimntei indiciando la fonte di provenienza della nuova licenz con l'avvertenza che essa contiene termini e condizioni differenti da quella originaria.

All'utente del software OS che intenda distribuire al pubblico il programma, \'e concessa anche la facolt\'a di richiedere un corrispettivo per le prestazioni di assistenza o di manutenzione che possono essere aggiunti quali ulteriori pattuizioni nella successiva licenza d'uso che accompagna il programma nella distribuzione. 

Per quanto attiene alle garanzie, i modelli di licenza OS sono generalmente privi di obblighi per il soggetto concendente ed il software viene concesso in unso nella stato cosiddetto "as it is".

A parte le caratteristiche comuni sopra delineate, le licenze OS presentano un contenuto variegato sia in termini di facolt\'a concesse al licenziatario che di limiti di utilizzo del software e della licenza stessa. 

A tal proposito esistono le licenze GPL e AGPL che rispettano la clausola copyleft, ed altre licenze come la Apache, MIT e BSD che sono dette invece permissive, nella quale non prevedono condizioni alla redistribuzion del software e permettendo di incorporare il codice in un software proprietario e rivenderlo al suo interno. 

Assai discussa \'e la natura giuridica delle licenze OS. Sembra prevalere la tesi per cui esse costituiscano un contratto gratuito atipico ni cui \'e riscontrabile un interesse economico del licenziante, che si realizza in modo indiretto attraverso la prospettiva di concudere altir negozi collegati, caratterizzati da un o scambio di prestazioni a carattere oneroso. 

Per quanto riguarda la commercializzazione dei software OS, esistono diversi modelli che vengono applicati dalle societ\'a di software. In alcuni casi, ad esempio, la societ\'a sceglie di distribuire e sviluppare gratuitamente il software OS, da combinarsi con altri programmi per elaboratore di tipo proprietario aventi funzionalit\'a maggiori. In altri casi, il produttore di software OS mira a ricavare degli utili dalle vendite al pubblico dei manuali esplicativi e della documentazione descrittivia inerente al programma. In altri casi, il produttore del software OS intende conseguire un risultato economico vantaggioso attraverso il servizio di personalizzazione del programma che pu\'o essere offerto a pagamento, o offrire all'utente (dietro corrispettivo) il servizio di assistenza e di manutenzione al programma per elaboratore. 

(fonte: https://assistenza-legale-imprese.it/software-tutela-giuridica-contratti-sviluppo-contratti-licenza/\#8.-Il-contratto-di-licenza-d-uso-del-software)

\subsection{Contratto d'opera: esempi}
Per foro competente si intende il Giudice competente a decidere della controversia. 

\subsubsection{Contratto d'opera: Articoli fondamentali}
Alcuni degli articoli fondamentali di un contratto d'opera di sviluppo software:

\begin{itemize}
    \item Sezione iniziale con eventuali premesse
    \item Articolo 1 - Oggetto del contratto: Il Fornitore si impegna a sviluppare un software secondo le specifiche tecniche dettate dal committente
    \item Articolo 2 - Durata e tempistiche: il fornitore si impegna a rispettare le seguenti scadenze: Analisi dei requisiti [data], Sviluppo [data inizio e fine], Testing e correzioni [data], Consegna finale [data]
    \item Articolo 3 - Corrispettivo e modalit\'a di pagamento: Pagamento con importo al momento della sottoscrizione del contratto, pagamento alla consegna della versione beta, saldo finale alla consegna ed accettazinoe del software
    \item Un'ulteriore sezione potrebbe essere aggiunta per la Penale: per ogni giorno di ritardo il Committente potr\'a richiedere o trattenere una determinata somma in euro.
    \item Modalit\'a di esecuzione dell'opera
    \item Modifiche in corso d'opera
    \item Accettazione dell'opera e difformit\'a
    \item Responsabilit\'a verso i terzi: Il prestatore d'opera \'e responsabile dei danni e/o pregiudizi arrecati al Committente o a terzi nell'esecuzione dei lavori
    \item Articolo 4 - Diritti di propriet\'a intellettuale: il Fornitore trasferice al Committente tutti i diritti di propriet\'a intellettuale sul software sviluppato, salvo eventuali componenti preesistenti o librerie di terze parti
    \item Articolo 5 - Obblighi delle parti: Il fornitore si impegna a rispettare le specifiche e le tempistiche concordate, garantire la qualit\'a del software, correggere eventuali difetti segnalati dal Committente entro un periodo di [numero] giorni dalla consegna
    \item Articolo 5.1 Il committente si impegna a: fornire tutte le informazioni ed il supporto necessari allo sviluppo del software, effettuare i pagamenti seconodo le modalit\'a stabilite
    \item Articolo 6: Riservatezza - Le parti si impegnao a mantenere riservate tutte le informazioni scambiate nell'ambito del presente contratto. La presente clausola rimarr\'a valida anche dopo la cessazione del contratto
    \item Articolo 7: Risoluzione del contratto: in quali casi il fornitore o il committente possono decidere di recedere dal contratto
    \item Sospensione dei lavori e impossibilit\'a sopravvenuta dell'esecuzione dell'opera
    \item Divieto di cessione del contratto: non subappalto
    \item Propriet\'a della documentazione e segreto professionale
    \item Risoluzione delle controversie
    \item Rinvio
    \item Spese contrattuali
    \item Tutela della Privacy
    \item Articolo 8: Legge applicabile e foro competente: Il contratto \'e regolato dalla legge italiano ed in caso di controversia dal foro compentente del luogo
    \item I diversi allegati per le specifiche tecniche, economiche, ecc.
    \item Firme
\end{itemize}

Un altro esempio di contratto potrebbe essere il seguente:
\begin{itemize}
    \item Clausola 1: Oggetto del contratto
    \item Clausola 2: Responsabilit\'a del fornitore
    \item Clausola 3: Rilascio versioni intermedie
    \item Clausola 4: Variazioni dei tempi previsti nel Piano di Lavoro
    \item Clausola 5: Verifica finale
    \item Clausola 6: Consegna
    \item Clausola 7: Titolarit\'a del codice sorgente (Ipotesi A, B, C, C1, D riportate sotto)
    \item Clausola 8: Formazione
    \item Clausola 9: Garanzia (manutenzione)
    \item Clausola 10: Revisioni (e patch) del software
\end{itemize}

Con il contratto d'opera una persona si obbliga a compiere, verso corrispettivo, un'opera o un servizio, con lavoro prevalentemente proprio e senza vincolo di subordinazione nei confronti del committente (art. 2222 c.c.). Nel contratto d'opera il soggetto si obbliga a pervenire ad un certo risultato, cno una conseguente assunzione del rischio nell'attivit\'a produttiva, rischio non conosciuto invece dal lavoratore dipendente.

L'ggetto del contratto non \'e il lavoro del prestatore d'opera, MA IL RISULTATO DELLO STESSO, a prescindere dalle modalit\'a sclete dall'obbligato per realizzarlo purch\'e l'opera corrisponda alle condizioni stabilite dal contratto e si sia proceduto a regola d'arte. 

Se l'esecuzione non \'e conforme a tali condizioni, il committente potr\'a avvalersi di un rimedio analogo alla diffida da adempiere ex art. 1454 c.c., e cio\'e la fissazione nei cnofronti del prestaotre d'opera di un congruo termine per adeguarsi alle stesse, con facolt\'a per il committente, in caso di decorso del termine, di recedere dal contratto, salvo in ogni cao il risarcimento del danno (art. 2224 c.c.).

Il legislatore, all'art. 2236 c.c., prevede poi una limitazione della responsabilit\'a del prestatore d'opere alle sole ipotesi di dolo o colpa grave nei casi in cui la prestazione implichi "la soluzione di problemi tecnici di speciale difficolt\'a ".

\subsubsection{Art. 2236 Codice Civile}
Se la prestazinoe implica la soluzione di problemi tecnici di speciale difficolt\'a, il prestatore d'opera non risponde dei danni, se non in caso di dolo o colpa grave, regolata dagli articoli 1176 e 2104 del c.c.

Per chiarire: con dolo, nel diritto civile, si intende un comportamento intenzionale volto a danneggiare, ingannare o frodare un'altra persona. Si distingue dalla colpa, che invece riguarda una condatta negligente o imprudente.

Per colpa grava, invece si intende un colpa particolarmente grave, caratterizzata da una violazione molto marcata del dovere di diligenza. Si tratta di un comportamento che, pur non essendo intenzinoale come il dolo, dimostra una trascuratezza spiccata e un evidente discostamento dalle regole di prudenza.

In sintesi: il dolo \'e una condotta intenzionale di danno, mentre la colpa grave \'e una trascuratezza spiccata. Entrambe possono portare alla responsabilit\'a civile, ma con alcune differenze a seconda del contesto e delle norme applicabili.

Anche se il prestatore d'opera non risponde dei danni, in casi di problemi tecnici di speciale difficolt\'a, deve comunque provvedere ad una diligenza qualificata, superiore a quella che viene richiesta ad una persona comune (c.d. diligenza del buon padre di famiglia), ed \'e commisurata alla prestazione che lo stesso deve eseguire. 

Per diligenza del buon padre di famiglia si intende l'impegno nel soddisfacimento dell'interesse del creditore, tipica dell'uomo 'medio', il buon padre di famiglia appunto, che va valutata in relazione alla specifica obbligazione che il debitore deve eseguire. La diligenza \'e comunque un concetto diverso da correttezza o buona fede (art 1175 c.c.). Queste ultime impongono alle parti di tenere un comportamente corretto enll'eseguire la propria prestazione ma non riguardano interessi specificamente predeterminati bens\'e il rapporto obbligatorio nel suo complesso. La diligenza, invece, indica le modalit\'a di esecuzione della prestazinoe e impone al debitore di fare tutto quanto necessario a soddisfare l'interesse del creditrore all'esatto adempimento. 

Nell'adempimento delle obbligazioni inerenti all'esercizio di un'attivit\'a professionale, la diligenza deve valutarsi con riguardo alla natura dell'attivit\'a esercitata. 

Il professionista risponde comunque di regola per negligenza, imprudenza e colpa lieve, atteso il maggior grado professionale che si presume in capo allo stesso. In ogni caso, la limitazione di responsabilit\'a non \'e applicabile alle ipotesi di negligenza o imprudenza del professionista, attesa la diligenza professionale richiesta allo stesso a norma dell'art. 1176, comma secondo. 

Per negligenza si intende la mancata attenzione, per imprudenza si intende l'agire in modo non sicuro e per imperizia si intende agire in modo non competente. La colpa lieve \'e una forma di colpa meno grave caratterizzata da una mancanza di diligenza, prudenza e perizia.

Nel caso in cui un cliente intenda agire per ottenere il risarcimento ha l'onere di provare il danno subito, la colpa del prestatore d'opera intellettuale, nonch\'e il nesso di causalit\'a tra colpa e danno. 

L'art. 2104 prevede che il prestatore di lavoro deve usare la diligenza richiesta dalla natura della prestazione dovuta, dall'interesse dell'impresa e da quello superiore della produzione nazionale [1176]. Deve inoltre osservare le disposizioni per l'esecuzione e per la disciplina del lavoro impartite dall'imprenditore e dai collaboratori di questo dai quali gerarchicamente dipende [2086, 2090, 2094, 2106, 2236]

Secondo l'art. 2224 del c.c. se il prestatore d'opera non procede all'esecuzione dell'opera secondo le condizioni stabilite [1622] dal contratto e a regola d'arte, il committente pu\'o fissare un congruo termine, entro il quale il prestatore d'opera deve conformarsi a tali condizioni [1454]. Trascorso inutilmente il termine fissato, il committente pu\'o recedere dal contratto [1373, 2227], salvo il diritto al risarcimento dei danni [1218]

Art. 2232 c.c. Il prestatore d'opera deve eseguire personalmente l'incarico assunto [1176]. Pu\'o tuttavia valersi, sotto la propria direzione e responsabilit\'a [1228], di sostituti e ausiliari (sempre e solo per l'esecuzione della prestazione affidatagli e rientrante nelle sue competenze), se la collaborazione di altri \'e consentita dal contratto o dagli usi e non \'e incompatibile con l'oggetto della prestazione [1717] (L'utilizzazione dei collaboratori ha valore interno: nei rapporti con i clienti rimane responsabile a tutti gli effetti solo il prestatore d'opera intellettuale).

Per regola d'arte si intende quell'insieme di tecniche dettate per l'esecuzione di un determinato genere di lavoro, considerate dai pi\'u corrette. Con "arte", si intende proprio la categoria professionale a cui appartiene il soggetto che esegue la lavorazione. Un'opera eseguita a regola d'arte \'e tale se il risultato prodotto \' e un manufatto sicuro, utilizzabile secondo l'uso programmato; per alcuni beni, \'e necessario anche il rispetto di determinati principi estetici e formali. Nell'esecuzione di opere professionali o artigianali, la regola d'arte va sempre rispettata quale dovere di diligenza ai sensi dell'art. 1176 del c.c.

Art.1454 c.c. Alla parte inadempiente l'altra pu\'o intimare per iscritto di adempiere in un congruo termine, con dichiarazione che, decorso inutilmente detto termine, il contratto s'intender\'a senz'altro risoluto. Il termine non pu\'o essere inferiore a quindici giorni, salvo diversa pattuizione delle parti o salvo che, per la natura del contratto o secondo gli usi, risulti congruo un termine minore. Decorso il termine senza che il contratto sia stato adempiuto, questo \'e risoluto di diritto.

La Diffida ad adempiere \'e quell'atto unilaterale e recettizio con cui la parte non inadempiente intima a quella inadempiente di eseguire la prestazione dovuta entro un termine congruo. Affinch\'e possa prodursi l'effetto di cui all'art. 1454 del c.c. (risoluzione del contratto) \'e necessario intimare l'adempimento avvertendo che, in caso contrario, il contratto si intender\'a risolto.

Art. 1373 c.c. Se una delle parti \'e attribuita la facolt\'a di recedere dal contratto, tale facolt\'a pu\'o essere esercitata finch\'e il contratto non abbia avuto un principio di esecuzione. Nei contratti a esecuzione continuata o periodica, tale facolt\'a pu\'o essere esercitata anche successivamente, ma il recesso non ha effetto per le prestazioni gi\'a eseguite o in corso di esecuzione. Qualora sia stata stipulata la prestazione di un corrispettivo per il recesso, questo ha effetto quando la prestazione \'e eseguita. \'E salvo in ogni caso il patto contrario. 

Art. 2227 c.c. Il committente pu\'o recedere dal contratto [1373, 2224], anchorch\'e sia iniziata l'esecuzione dell'opera, tenendo indenne il prestore d'opera delle spese, del lavoro eseguito e del mancato guadagno.

Art. 1372 c.c. Il contratto ha forza di legge tra le parti. Non pu\'o essere sciolto che per mutuo consenso o per cause ammesse dalla legge. Il contratto non produce effetto rispetto ai terzi che nei casi previsti dalla legge []

Art. 1671 c.c. Il committente pu\'o recedere dal contratto, anche se \'e stata iniziata l'esecuzione dell'opera o la prestazione del servizio, purch\'e tenga indenne l'appaltatore delle spese sostenute, dei lavori eseguiti e del mancato guadagno. La norma non pone alcuna condizione per l'esecizio del recesso che, quindi, \'e libero.

Con il termine "indenne" si intende l'indennizzo che il committente deve versare nei confronti del prestatore d'opera. Con "Indennit\'a (o indennizzo)" si intende la somma di denaro dovuta ad un soggetto per un pregiudizio da lui subito che non consegue ad un atto illecito e quindi a responsabilità civile, ma viene conseguita a titolo di ristoro patrimoniale che consegue a fatti che sacrificano diritti altrui ma che non sono antigiuridici in quanto autorizzati o imposti da una norma di legge. Pertanto, va tenuta distinta dal risarcimento dei danni sia in quanto non consegue alla violazione di un obbligo, sia in quanto non è diretta a reintegrare in pieno la sfera giuridica sacrificata.

Art 1734 c.c. Il committente pu\'o revocare l'ordine di concludere l'affare fino a che il commissionario non l'abbia concluso. In tal caso spetta al commissionario una parte della provvigione, che si determina tenendo conto delle spese sostenute e dall'opera prestata.

Art 2237 c.c. Il cliente pu\'o recedere dal contratto, rimborsando al prestaotre d'opera le spese sostenute e pagando il compenso per l'opera svolta. Il prestatore d'opera pu\'o recedere dal contratto per giusta causa [2119]. In tal caso egli ha diritto al rimborso delle spese fatte e al compenso per l'opera svolta, da determinarsi con riguardo al risultato utile che ne sia derivato al cliente. Il recesso del prestatore d'opera deve essere esercitato in modo da evitare pregiudizio al cliente.

Art 2119 c.c. Ciascuno dei contraenti pu\'o recedere dal contratto [1373] prima della scadenza del termine, se il contratto \'e a tempo determinato [2097], o senza preavviso, se il creottao \'e a tempo indeterminato, qualora si verifichi una causa che non consenta la prosecuzione, anche provvisoria, del rapporto [2103, 2244]. Se il contratto \'e a tempo indeterminato, al prestore di lavoro che recede per giusta causa compete l'indennit\'a indicata nel secondo comma dell'articolo precedente. Non costituisce giusta causa di risoluzione del contratto la liquidazione coatta amministrativa dell'impresa. Gli effetti della liquidazione giudiziale sui rapporti di lavoro sono regolati dal Codice della crisi d'impresa e dell'insolvenza. 

Art. 1672 c.c. Se il contratto si scioglie perché l'esecuzione dell'opera è divenuta impossibile in conseguenza di una causa non imputabile ad alcuna delle parti, il committente deve pagare la parte dell'opera già compiuta, nei limiti in cui è per lui utile [1675, 2228], in proporzione del prezzo pattuito per l'opera intera (della parte di opera di cui paga il prezzo il committente diviene, ovviamente, proprietario).

Art. 2228 c.c. Se l'esecuzione dell'opera diventa impossibile (impossibilit\'a oggettiva quando riguarda una delle parti del contratto e determina l'estinzione perch\'e \'e determinante la persona; oggettiva quando rende impossibile la prosecuzione dell'opera) per causa non imputabile ad alcuna delle parti [1464](2), il prestatore d'opera ha diritto ad un compenso per il lavoro prestato in relazione all'utilità della parte dell'opera compiuta [1672, 2231, 2237].

Art. 1218 c.c. (1)Il debitore che non esegue esattamente(2) la prestazione dovuta [1176, 1181] è tenuto al risarcimento del danno [1223 ss.], se non prova(3) che l'inadempimento o il ritardo(4) è stato determinato da impossibilità della prestazione derivante da causa a lui non imputabile [1221, 1229, 1257, 1307, 1557, 1558, 1673, 1693, 1821, 2740; 160 disp. trans.](5)(6).

Art. 1776 c.c. L'erede del depositario, il quale ha alienato in buona fede la cosa che ignorava essere tenuta in deposito, è obbligato soltanto a restituire il corrispettivo ricevuto [1811]. Se questo non è stato ancora pagato, il depositante subentra nel diritto dell'alienante [1203]. 

Art 1228 c.c. Salva diversa volontà delle parti, il debitore che nell'adempimento dell'obbligazione si vale dell'opera di terzi, risponde anche dei fatti dolosi o colposi di costoro [1229, 1717, 2049].

[2222] Contratto d'opera
[2236] Dolo o colpa grave
[1176] Diligenza del buon padre di famiglia
[2104] Diligenza
[2224] Termine di conformazione
[2232] Lavoro proprio
[1454] risoluzione del contratto
[1372] Il contratto ha forza di legge tra le parti
[1373] Recessione da contratto
[1671] Recesso unilaterale dal contratto
[1734] Revoca dell'affare
[2237] Recesso dal contratto
[2227] Recesso unilaterale dal contratto (per prestatore d'opera)
[2231] Recesso per mancata iscrizione all'albo (da approfondire)
[2119] Recesso per giusta causa
[1672] Risoluzione contratto per impossibilit\'a di prosecuzione
[2228] Impossibilit\'a di esecuzione dell'opera
[1218] Risarcimento del danno per inadempimento o ritardo (da approfondire)
[1776] 
[1228] Responsibilit\'a di terzi
[1717]
[2103]
[2244]
[2097]
[1672]
[2231]
[2237]
[1229]
[1717]
[2049]
[1203]
[1811]
[1221]
[1229]
[1257]
[1307] 
[1557]
[1558]
[1673]
[1693]
[1821]
[2740]
[160 disp. trans.]

\subsubsection{Contratto d'appalto: esempi}
\begin{itemize}
    \item Articolo 1: Premesse
    \item Articolo 2: Oggetto dell'appalto
    \item Articolo 3: Ammontare dell'appalto
    \item Articolo 4: Modalit\'a di liquidazione
    \item Articolo 5: Tracciabilit\'a dei flussi finanziari e modalit\'a di pagamento
    \item Articolo 6: Termini di esecuzione e penali
    \item Articolo 7: Garanzia definitiva
    \item Articolo 8: Cessione del contratto e cessinoe dei crediti
    \item Articolo 9: Periodo di assistenza tecnica e manutenzione
    \item Articolo 10: Formazione del personale
    \item Articolo 11: Responsabile di contratto per l'appaltatore
    \item Articolo 12: Codice di comportamento
    \item Articolo 13: Subappalto
    \item Articolo 14: Risoluzione del Contratto
    \item Articolo 15: Recesso dal contratto
    \item Articolo 16: Elezione di domicilio
    \item Articolo 17: Controversie
    \item Articolo 18: Obbligo di segretezza e riservatezza
    \item Articolo 19: Trattamento dati personali
    \item Articolo 20: Spese, imposte e tasse
    \item Articolo 21: Patto di integrit\'a
    \item Firme Appaltatore e Appaltante
\end{itemize}

Art. 1667 c.c. L'appaltatore \'e tenuto alla garanzia per le difformit\'a e i vizi dell'opera. La garanzia non \'e dovuta se il committente ha accetta l'opera [1665] e le difformit\'a o i vizi erano da lui conosciuti o erano riconoscibili, purch\'e in questo caso, non siano stati in malafede taciuti dall'appaltatore.

Il committente deve, a pena di decadenza, denunziare all'appaltatore le difformit\'a o i vizi entro sessanta giorni dalla scoperta. La denunzia non \'e necessaria se l'appaltatore ha riconosciuto le difformit\'a o i vizi o se li ha occultati [1670]. 

L'azione contro l'appaltatore si prescrive in due anni dal giorno della consegna dell'opera. Il comittente convenuto per il pagamento pu\'o sempre far valere la garanzia, purch\'e le difformit\'a o i vizi siano stati denunciati entro sessanta giorni dalla scoperta e prima che siano decorsi i due anni dalla consegna. [181 disp. att.; 240,  855 c. nav.]

Con difformit\'a si indica la rituazinoe in cui l'opera ultimata non corrisponde esattamente, dal punto di vista tecnico, all'opera oggetto del contratto. L'esistenza di difformit\'a fa s\'i che il committente possa invocare la garanzia di cui all'art. 1667 del c.c., in quanto anche le prescrizioni tecniche (solitamente contenute in un progetto) fanno parte integrante del contenuto contrattuale e quindi devono essere rispettate dall'appaltatore.

In un contratto d'appalto, il vizio sorge in fase di esecuzione dell'opera o del manufatto, principlametne in quanto il costrutto re viola le regole dell'arte o i principi in materia di edilizia. Il committente deve essere garantito dall'appaltatore contro i vizi occulti e non facilmente riconoscibili, avendo il diritto di chiedere che le difformit\'a o i vizi siano eliminati a sese dell'appaltatore, oppure che il prezzo sia proporzionalmente diminuito, salvo il risarcimento del danno nel caso di colpa dell'appaltatore. Se le difformit\'a o i vizi dell'opera sono tali da renderla del tutto inadatta alla sua destinazione, il committente pu\'o addirittura chiedere la risoluzione del contratto. Egli \'e per\'o tenuto a fare denunzia dei vizi entro un 60 giorni dalla scoperta e deve agire a pena di prescrizione entro i successivi due anni.

\subsection{Contratto di Licenza d'Uso}
Elementi chiave di un contratto di licenza d'uso di sofware con canone periodico:
\begin{itemize}
    \item Parti: Licenziante e Licenziatario
    \item Oggetto: diritto di utilizzare il software
    \item Canone: Pagamento ricorrente per l'uso del software
    \item Durata: Il contratto pu\'o essere a tempo determinato o indeterminato, con clausole di rinnovo o recesso
    \item Limitazinoi: le condizioni di utilizzo del software (es. numero di utenti, tipi di dispositivi, ecc.) sono specificate nel contratto
    \item Obblighi del licenziatario: Pagamento del canone, rispetto delle condizioni di utilizzo, e manteninmento della riservatezza
    \item Obblighi del licenziante: Concessione del diritto di utilizzo, eventuali servizi di supporto o aggiornamenti
    \item Responsabilit\'a: Le parti possono essere responsabili per eventuali violazioni del contratto
\end{itemize}

Un esempio concreto di contratto:
\begin{itemize}
    \item Definizioni
    \item Concessione della Licenza
    \item Limitazioni
    \item Informazioni sul Copyright
    \item Propriet\'a
    \item Limitazioni relative all'esportazione
    \item Supporto e Aggiornamento
    \item Esclusione di garanzia e Responsabilit\'a
    \item GARANZIA LIMITATA
    \item NESSUNA RESPONSABILIT\'A IN CASO DI DANNI
    \item LIBERATORIA DI RESPONSABILIT\'A
    \item Cessazione del Contratto
    \item AVVISO DI DIRITTI LIMITATI DEL GOVERNO USA
    \item Disgiungibilit\'a
    \item Software di terzi
    \item Accettazione
\end{itemize}

Un altro esempio di contratto di licenza d'uso
\begin{itemize}
    \item Le parti: licenziante e licenziatario
    \item Premesse
    \item Definizioni
    \item Oggetto del contratto
    \item Propriet\'a intellettuale e limiti della licenza d'uso
    \item Utilizo dei modelli software e titolarit\'a degli stessi
    \item Durata del contratto
    \item Modifica, riproduzione e utilizzo
    \item Garanzia limitata
    \item Limitazione di responsabilit\'a ed esclusione danni
    \item Corrispettivo, fatturazione e pagamento
    \item Attivazinoe della licenza
    \item Trattamento dei dati
    \item Utilizzo dati non personali
    \item Legge applicabile e Foro competente
    \item Divieto di cessione
    \item Clausola risolutiva espressa
    \item Allegati
\end{itemize}

I diversi tipi di licenza sono:
\begin{itemize}
    \item Licenza perpetua: In questo tipo di contratto il licenziatario ha il diritto di utilizzare il software senza limiti di tempo. Questo tipo di licenza \'e comune per software che non richiedono aggiornamenti continui
    \item Licenza temporanea: il licenziante concende il diritto di utilizzo del software per un determinato periodo di tempo, oltre il quale la licenza dovr\'a essere rinnovata.
    \item Licenza basata su abbonamento: in questo caso il licenziatario paga un corrispettivo in denario con scadenze regolari per mantenere l'accesso al software.
\end{itemize}

\subsubsection{Licenza d'Uso: Software as a Service}
\subsubsection{Licenza d'Uso: licenza a strappo}
\subsubsection{Licenza d'Uso: licenza OEM}

\subsection{Licenza Open Source}

\subsection{Inizio Parte da sistemare}


\subsection{Propriet\'a intellettuale}
legge sul diritto d'autore
riproduzione

codice oggetto e codice sorgente
Alla software house spettano sempre i diritti morali, che non possono essere ceduti.

I diritti di sfruttamenteo economico, invece, possono essere ceduti tramite contratto, tramite royalty ad esempio.

Spesso quando si acquista un software, se ne acquista il diritto d'uso, non del codice in s\'e, sia quello oggetto che sorgente.

Anche nel caso in cui si sviluppi il codice per un cliente terzo, \'e necessario stabilire chi ha il possesso del codice sorgente, il cliente o lo sviluppatore?

In determinati casi, \'e anche possibile brevettare il codice sorgente.

\subsection{Contratto di licenza d'uso}
diritto di utilizzo del programma a fronte di un corrispettivo
dottrina minoritaria: innominato
dottrina prevalente: contratto di locazione

Licenziante e Licenziatario

Senza trasferire la propriet\'a del software.

accordi licenziante al licenziatario, senza trafroirne la propriet\'a

A cosa servono? consentire a terzi di usare il software, senza cedere la propriet\'a

commercializzati per uso da parte di terzi

scalabile, controllo sul software

canone per il solo uso del software

per sempre o per un certo periodo di tempo
in acluni casi, accesso al codice sorgente

contratti di compravendita del software. 

contratti di sviluppo di software

qauli diriti spettino alle parti

brevetti, non sono considerati invenzioni
software non brevettabile in se 



\subsection{Contratto di Compravendita}
contratto di compravendita
oggetto d'acquisto del bene immateriale
pieno e illimitato potere sul programma
singola copia del software
pu\'o rivenderlo a terzi
vendita tecnica

\subsection{Contratto di Sviluppo}
\subsubsection{Contratto d'opera}
Contratto d'opera
contratto di sviluppo: il programmatore elabori il software su richieste
se si tratta di un professionista: contratto d'opera
professionista sviluppare il progrmama con lavoro proprio
mezzi e responsibilita dolo o colpa grave

\subsubsection{Contratto d'appalto}
Contratto d'appalto
gestione a proprio rischio
per vizi la questione \'e ben pi\'u complessa
gestione a prorprio rischio
conformita rispettoa lle richieste tecniche

\subsection{Brevetti}

\subsection{Contratto di licenza d'uso software (EULA - End User License Agreement)}
\subsection{Contratto di servizio SaaS (Software as a Service)}
\subsection{Contratto di manutenzione e assistenza}
\subsection{Contratto di personalizzazione}
\subsection{Privacy Policy e trattamento dati}

\subsection{Contratto di sviluppo software / realizzazione app}
\subsection{Contratto di consulenza informatica}
\subsection{Contratto di manutenzione e supporto tecnico}
\subsection{Accordo di riservatezza (NDA)}


\section{Cassa previdenziale}
\subsection{Contributi previdenziali}
\subsubsection{Contributi con Cassa Previdenziale}

\subsubsection{Contributi con Gestione Artigiani e Commercianti INPS}
L'attivit\'a di artigiano o commerciante richiede l'iscrizione alla Camera di Commercio e il versamento dei contributi alla Gestione Artigiani e Commercianti dell'INPS. 

Questa cassa prevede contributi fissi da pagare annualmente indipendetemente dal reddito percepito.

Per l'INPS il reddito minimale per il 2025 sui cui si pagano i contributi \'e di 18.555 euro. 

Per artigiani in regime forfettario il contributo fisso nel 2025 \'e di 4.460,64 euro (4.371.80 euro per chi ha meno di 21 anni).

Per i commercianti in regime forfettario: 4.549,70 euro (4.460,19 euro per chi ha meno di 21 anni).

Per il reddito eccedente il reddito minimale (di 18.555) \'e prevista un'aliquota aggiuntiva,

Fino al reddito di 55.448 euro le aliquote sono: 24,48\% (24,18\% per chi ha meno di 21 anni) per i commercianti e 24\% (23,70\% per chi ha meno di 21 anni) per artigiani. 

Oltre i 55.448 euro le aliquote sono 25,48\% (25,18\% per chi ha meno di 21 anni) per i commercianti e 25\% (24,70\% per chi ha meno di 21 anni) per gli artigiani. 

Esclusivamente per artigiani e commercianti che operano in regime forfettario \'e prevista la possibilit\'a di riduzione dei contributi INPS. Esistono due tipi di riduzione, una del 35\% ed un'altra del 50\% dei contributi previdenziali.

La riduzione del 35\% non ha limiti di tempo, ma si pu\'o fare richiesta una sola volta, quindi se si rinuncia a questa agevolazione, non si pu\'o pi\'u cambiare idea. Per aderire bisogna esplicitamente fare richiesta all'INPS. Questa riduzione si applica sia ai contributi fissi che variabili.

Inoltre per chi si iscrive per la prima volta in Camera di Commercio, come artigiano o commerciante, \'e prevista una riduzione del 50\% dei contributi per i primi 36 mesi di attivit\'a.

Per chi non ha ancora una partita iva, deve presentare domanda dopo il 28 febbraio, immediatamente dopo aver ricevuto il provvedimento di apertura della posizione previdenziale alla gestione Artigiani e Commercianti dell'INPS.

Il calcolo da effettuare per la determinazione del contributo finale \'e la seguente: ( (contributi fissi) + (contributi eccedenti) ) * (percentuale di riduzione). 

L'unico inconveniente per chi aderisce a questo tipo di riduzione \'e il fatto che un anno di contributi a livello previdenziale potrebbe non corrispondere ad un anno effettivo, proprio per via della riduzione.

Supponendo che un artigiano non superi la soglia di reddito minale, allora si ritrover\'a a dover pagare contributi per: 
\begin{itemize}
    \item 4.460,64 euro * 35\% = 3.016,38 euro
    \item 4.640,64 euro * 50\% = 2.320,32 euro
\end{itemize}

\subsubsection{Contributi con Gestione separata INPS}

\subsection{Contributi INPS}
L’INPS gestisce la liquidazione e il pagamento delle pensioni e delle indennità di natura previdenziale e assistenziale. Le pensioni sono prestazioni previdenziali, determinate sulla base di rapporti assicurativi e finanziate con i contributi di lavoratori e aziende pubbliche e private. Invece, le prestazioni assistenziali o “a sostegno del reddito” tutelano i lavoratori che si trovano in particolari momenti di difficoltà della loro vita lavorativa e provvedono al pagamento di somme destinate a coloro che hanno redditi modesti e famiglie numerose. Per alcune di queste prestazioni l’INPS è coinvolto solo nella fase di erogazione, mentre per altre svolge tutto il procedimento di assegnazione (fonte: sito INPS). L’INPS nasce oltre centoventicinque anni fa, nel 1898 con la fondazione della Cassa Nazionale di previdenza per l'invalidità e per la vecchiaia degli operai, allo scopo di garantire i lavoratori dai rischi di invalidità, vecchiaia e morte. Con il tempo l’Istituto ha assunto un ruolo di crescente importanza fino a diventare il pilastro del sistema nazionale del welfare.

\subsection{Contributi INAIL}
L'INAIL (Istituto Nazionale per l'Assicurazione contro gli Infortuni sul Lavoro) \'e un ente pubblico italiano che si occupa di assicurazione contro gli infortuni sul lavoro.
In generale offre servizi per la tutela dei lavoratori in caso di danni alla salute causati dallo svolgimento dell'attivit\'a lavorativa, promuove iniziative per migliorare la sicurezza nei luoghi di lavoro, 
fornisce supporto medico, riabilitativo ed anche professionale ai lavoratori infortunati o malati ed eroaga compensazioni economiche in caso di infortunio o malattia professinoale che comporti inabilit\'a temporanea, permanente o, nei casi pi\'u gravi, la morte del lavoratore.

Per alcune attivit\'a considerate rischiose (artigianato, edilizia, industria, ecc.), l'assicurazione INAIL \'e obbligatoria (da parte dei datori di lavoro).

\section{Documenti}

\subsection{Modello F24}
Il modello F24 \'e il modulo ufficiale usato in Italia per pagare tutte le imposte, tasse e contributi dovuti allo Stato, all'INPS, alle Regioni e ad altri enti pubblici.

Serve per pagare IRPEF, Addizionale comunale IRPEF, Imposta comunale sugli immobili (ICI), IMU, Tarsu, Tosap/cosap, Interessi pagamenti rateali, Addizionale regionale IRPEF, Tares, Tasi.

In maniera molto riassuntiva \'e composto da diverse sezioni, quali i dati anagrafici del contribuente, i codici tributo, gli importi a debito o a credito ed il periodo di riferimento.

Il modello F24 pu\'o essere pagato tramite home banking o Fisconline, tramite commercialista oppure in banca o alla posta.

Esistono diversi tipi di modello F24 tra cui ordinario utilizzato per la maggior parte dei pagamenti, il semplificato utilizzato per tributi locali come IMU, TARI, ecc. ed il precompilato fornito solitamente dal commercialista.

Per il pagamento di un determinato tributo, almeno in maniera minima, sono necessari i dati del contribuente (codice fiscale, nome, cognome, indirizzo, comune, provincia, CAP) 
ed i dati della sezione INPS (codice ente, causale contributo, codice tributo, anno di imposta, importo a debito).   

Il modello F24 \'e un modulo utilizzato per il pagamento di tasse, imposte, contributi e tributi e serve a versare importi dovuti all'Agenzia delle Entrate, all'INPS, ai Comuni o ad altri enti. 

Attraverso il modello F24 si possono pagare: imposte sui redditi (IRPEF, IRES), IVA, IMU, TASI, TARI, Contributi INPS e INAUL, Sanzioni e Interessi, Accise e tributi regionali e locali

Esistono diverse versioni del modello F24: F24 ordinario, F24 semplificato, F24 accise ed F24 ELIDE.

Il modello F24 pu\'o essere pagato Online tramite i servizi telematici dell'Agenzia delle Entrate o del proprio home banking, in banca o posta o tramite intermediari abilitati come i commercialisti.


\subsubsection{Esempi di compilazione modello F24}
Per le imposte, supponendone il normale pagamento,
queste possono essere IRPEF saldo, IRPEF acconto prima rata, IRPEF acconto seconda rata o acconto in unica soluzione.

Per il pagamento IRPEF saldo bisogna compilare la sezione erario, ed inserire i seguenti dati: 
Codice Tributo(4001), 
rateazione (n.rata - numero di rate in cui si vuole frazionare l'importo totale, e.g. '0107'), 
anno di riferimento, 
importi a debito versati, 
importi a credito compensati (NON COMPILARE),
TOTALE A (Somma degli importi a debito indicati nella Sezione Erario),
TOTALE B (Somma degli import a credito indicati nella Sezione Erario, non compilare se non presenti importi a credito)
+/- (indicare il segno se positivo o negativo),
SALDO (A - B) (indicare il saldo TOTALE A - TOTALE B),
codice ufficio (NON COMPILARE),
codice atto (indicare il Codice dell'Atto oggetto di definizione)

Esempio imposta di bollo:

Esempi contributi:


\section{Consulente}
\subsection{Il commercialista}
Cosa fa: si occupa della tenuta della contabilit\'a, 
predisposizione del bilancio di esercizio e della dichiarazione dei redditi, 
redazione dei libri contabili e libro unico del lavoro, 
svolgimento di perizie tecniche, ispezioni e revisioni amministrative.

La professione del commercialista pu\'o essere definita come quella del consulente che opera 
nel campo del diritto commerciale, del diritto tributario, della ragioneria e della contabilit\'a in generale.

Il diritto commerciale \'e una branca del diritto privato che regola i rapporti attinenti alla produzione e allo scambio della ricchezza. Pi\'u in particolare, regola ed ha per oggetto i contratti conclusi tra operatori economici e tra essi ed i loro i clienti privati ed anche gli atti e le attivit\'a delle societ\'a. 

Il diritto tributario \'e un settore del diritto finanziario che regolamente i tributi di ogni tipo.

La ragioneria \'e una disciplina che si occupa della rilevazinoe dei fenomeni aziendali e della loro traduzione in scritture cnotabili, allo scopo di fornire alle direzioni delle aziende 
o ad altri soggetti interessati le informazioni necessarie per esercitare le funzioni di controllo e di decisione.

Per contabilit\'a si intende al storia economica di un'azienda nel senso che conserva una traccia di tutte le operazioni commerciali.

Oltre al commercialista, per consulenza fiscale \'e possbili avvalersi del CAF, dei consulenti del lavoro, o direttamente all'Agenzia delle Entrate.

I costi tipici di un commercialista variano in base alla forma societaria, al regime fiscale, al prestigio del commercialista stesso ed altri parametri come tipologia del servizio, numero di fatture, settore in cui svolgi l'attivit\'a ecc.

Una partita IVA in regime forfettario paga intorno ai 60 euro al mese,
un titolare di partita IVA in regime semplificato potrebbe pagare dai 900 ai 2000 euro annui, mentre
un titolare di societ\'a o lavoratore autonomo in regime ordinario tra i 2000 ed i 4000 euro annui (tra i 150 ed i 350 euro al mese).


\subsection{Il CAF}

\subsection{Agenzia delle entrate}



\section{Contenziosi}

\subsection{Contenziosi: Agenzia delle Entrate}
\begin{itemize}
    \item Omissioni o errori nella fatturazione
    \item Compensi non dichiarati
    \item Superamento soglie del regime forfettario, ma non si passa al regime ordinario
    \item Deducibilità di spese non inerenti (viaggi, cene, prodotti ad uso personale non legate all'attivit\'a produttiva)
    \item Mancata iscrizione alla gestione previdenziale INPS (o errata)
    \item Errori nei versamenti di imposte e contributi
    \item Errori nei rapporti con clienti esteri (mancata iscrizione al VIES, reverse charge non applicato)
\end{itemize}

\subsubsection{Cartelle esattoriali}
Le cartelle esattoriali sono atti ufficiali di riscossione meessa dall'agenzia delle Entrate per richiesere il pagamento di debiti verso lo Stato o altri enti pubblici. In genere contengono: l'importo da pagare (comprensivo di debito originario, sanzioni, interessi di mora, eventuali spese di notifica e riscossione), l'ente creditore (es. INPS, Comune, Agenzia delle Entrate, ...), le istruzioni per il pagamento, il termine per pagare o fare ricorso. 

Esistono diversi motivi per cui si potrebbero ricevere delle cartelle esattoriali, come ad esempio: Omissioni o ritardi nei versamenti fiscali (IRPED, IVA, imposta sostitutiva forfettaria), Contribuit INPS non versati, Sanzioni per errori fiscali, multe per mancata comunicazione (esterometro, CU omesse, etc.), debiti pregressi di cui si \'e a conoscenza. Se non si paga una determinata cartella esattoriale si rischia un fermo amministrativo su un auto, il pignoramento di un conto corrente o dello stipendio, o un'ipoteca su immobili. 

Ricevendo una cartella, \'e bene verificare che sia corretta, controllando che il debito sia erale e verificare che non sia caduto in prescrizione, si pu\'o rateizzare il pagamento o si pu\'o fare ricorso entro 60 giorni se la si ritiene errata. 


\subsection{Contenziosi: Clienti}
\begin{itemize}
    \item Mancato pagamento o ritardo nei pagamenti
    \item Ambiguità o assenza di contratto
    \item Scope creep (espansione non concordata delle richieste)
    \item Diritti di propriet\'a intellettuale
    \item Violazione della privacy o del GDPR
    \item Contenziosi con agenzie o intermediari
    \item Conflitti sulla qualit\'a del software
\end{itemize}

\end{document}